% !TEX root = ../../main.tex
% =====================================================================
% Chapter: The Shadow Tower -- Quadrichotomy
%
% Master narrative anchored by the Quadrichotomy Theorem and the
% Riccati master equation H^2 = t^4 Q_c. Consolidates the 15+
% scattered shadow_tower_*.tex files into one cathedral with a
% single load-bearing display, partition theorem, and five
% corollaries. Frontier inscriptions of the Borel-geometric layer
% (iter 30-77) cited as subordinated lemmata.
%
% Subordinated chapters (NOT modified, only oriented around):
%   shadow_tower_higher_coefficients.tex          (S_6, S_7, S_8 closed)
%   shadow_tower_sub_subleading_platonic.tex      (Gamma_r layer)
%   shadow_tower_other_class_M_platonic.tex       (W_3, W_N, W_inf, BP)
%   motivic_shadow_tower.tex                      (motivic lift)
%   motivic_shadow_full_class_m_platonic.tex      (E_1 depth-1 thm)
%   shadow_L_function_platonic.tex                (AP70 healed)
%   conformal_anomaly_rigidity_platonic.tex       (kappa(Vir)=c/2 rigidity)
%   chern_weil_level_shift_platonic.tex           (kappa to genus tower)
%   clutching_uniqueness_platonic.tex             (AP225)
%   z_g_kummer_bernoulli_platonic.tex             (Z_g closed forms)
%   higher_kummer_arithmetic_duality_platonic.tex (extended duality)
%   all_tier_generating_function_platonic.tex     (Tier-K closed form)
%   virasoro_motivic_purity_all_r_platonic.tex    (rationality of S_r)
%
% Standards followed: ClaimStatus tags, label discipline,
% kappa qualified per family, `\providecommand`
% for portability, no AI slop, no em dashes,
% shadow L-function poles respected, Borel-summability
% qualifier (convergent series uses Pade not Borel), on the
% S_4 large-c asymptotic, AP1 on never writing kappa from memory
% (each formula carries a landscape census citation comment).
%
% Author: Raeez Lorgat.
% =====================================================================

\providecommand{\ChapterIntro}[1]{\par\noindent\textit{#1}\par\medskip}

\providecommand{\cA}{\mathcal{A}}
\providecommand{\cF}{\mathcal{F}}
\providecommand{\cW}{\mathcal{W}}
\providecommand{\cC}{\mathcal{C}}
\providecommand{\cO}{\mathcal{O}}
\providecommand{\cM}{\mathcal{M}}
\providecommand{\fg}{\mathfrak{g}}
\providecommand{\fh}{\mathfrak{h}}
\providecommand{\fL}{\mathfrak{L}}
\providecommand{\Vir}{\mathrm{Vir}}
\providecommand{\Walg}{\mathcal{W}}
\providecommand{\BP}{\mathrm{BP}}
\providecommand{\Heis}{\mathrm{Heis}}
\providecommand{\KM}{\mathrm{KM}}
\providecommand{\hv}{h^{\vee}}
\providecommand{\MC}{\mathrm{MC}}
\providecommand{\Sh}{\mathrm{Sh}}
\providecommand{\wt}{\mathrm{wt}}
\providecommand{\Res}{\operatorname{Res}}
\providecommand{\disc}{\operatorname{disc}}
\providecommand{\str}{\operatorname{str}}
\providecommand{\Pic}{\mathrm{Pic}}
\providecommand{\PF}{\mathrm{PF}}
\providecommand{\Bor}{\mathrm{Bor}}
\providecommand{\ClaimStatusProvedHere}{}
\providecommand{\ClaimStatusProvedElsewhere}{}
\providecommand{\ClaimStatusConjectured}{}
\providecommand{\ClaimStatusHeuristic}{}
\providecommand{\ClaimStatusConditional}{}

\chapter{The Shadow Tower: Quadrichotomy}
\label{chap:shadow-quadrichotomy-platonic}

\index{shadow tower!quadrichotomy theorem|textbf}
\index{Riccati master equation|textbf}
\index{Picard--Fuchs!shadow connection|textbf}
\index{spectral curve!hyperelliptic Sigma_c|textbf}

\ChapterIntro{The shadow obstruction tower
$\{S_r(\cA)\}_{r\ge 2}$ of a chirally Koszul vertex algebra $\cA$,
restricted to a primary line $L$ with non-vanishing curvature
$\kappa = S_2(\cA|_L)$, is governed by a single algebraic identity:
the Riccati equation $H(t)^2 = t^4 Q_c(t)$, with
$Q_c(t) = (2\kappa + 3\alpha\,t)^2 + 2(8\kappa S_4)\,t^2$. The
right-hand side is determined by three numbers $(\kappa, \alpha, S_4)$
read off from the OPE at the level of the 2-, 3-, and 4-point
function. The four classes G/L/C/M of the standard landscape are the
four factorisation types of $Q_c$, in codimensions $2, 1, 1, 0$ of
the parameter space $\mathbb A^3 = \{(\kappa,\alpha,S_4)\}$. Class M
acquires a hyperelliptic spectral Riemann surface
$\Sigma_c = \{y^2 = Q_c\}$ whose periods are the flat sections of the
shadow connection $\nabla^{\rm sh}$. Five corollaries follow from the
master equation: $\Theta_\cA$ is a Maurer--Cartan element; $\Sigma_c$
governs the genus tower via Picard--Fuchs; the asymptotic series is
Gevrey-$1$ and admits a directional resummation by Pad\'e or Borel
where applicable; the shadow--Feynman dictionary
$F_L = \mathrm{Sh}_{L+1}$ is a Getzler--Kapranov boundary identity;
the $\cW_N$ characters at admissible levels are conjecturally mock
modular by Borcherds lift on $\Sigma_c$. The Borel-geometric layer
(iter~$30$--$77$) provides the explicit constants:
$C_\cA = 6$ universal on the non-logarithmic class M $T$-line,
$C_{\cW_3}^{W{\rm-line}} = 12 = 2\cdot 6$ on the second lane,
$\beta_N = 12(H_N - 1)$ per-spin lane in $\cW_N$.}

\section{The shadow tower as MC obstruction}
\label{sec:shadow-mc-obstruction}

The shadow obstruction tower $\{S_r(\cA)\}_{r\ge 2}$ is the sequence
of Maurer--Cartan coefficients projected from the bar-intrinsic
element $\Theta_\cA$ to a primary line $L$.

\begin{definition}[Shadow tower on a primary line]
\label{def:shadow-tower-primary-line}
Let $\cA$ be a chirally Koszul vertex algebra of finite type, and let
$L \subset \cA$ be a primary line equipped with a generator $x$ of
conformal weight $\Delta$. The
\emph{shadow tower} $\{S_r(\cA, L)\}_{r\ge 2}$ on $L$ is the sequence
of MC-coefficients defined by
\[
\Theta_{\cA}\big|_{L} \;=\; \sum_{r\ge 2} S_r(\cA, L)\,x^{r}
\;\in\; \mathfrak g^{\rm mod, (0)}_{\cA}\big|_{L},
\]
where $\Theta_{\cA}$ is the bar-intrinsic Maurer--Cartan element
(Theorem~\ref{thm:mc2-bar-intrinsic} of the bar-complex part). When
the line $L$ is fixed by context we write $S_r(\cA)$.
\end{definition}

\begin{remark}[Curvature]
\label{rem:curvature-ap1-ap39}
The leading coefficient $\kappa(\cA, L) := S_2(\cA, L)$ is the
modular characteristic on $L$. For Virasoro on its single primary
line we have
\[
\kappa(\Vir_c) \;=\; \tfrac{c}{2}
\qquad \text{(landscape census C3 / kappa(Vir) line).}
\]
For the affine Kac--Moody algebra $V_k(\fg)$ on the universal vacuum
line,
\[
\kappa(V_k(\fg)) \;=\; \frac{\dim(\fg)\,(k + \hv)}{2\,\hv}
\qquad \text{(landscape census C3).}
\]
At $k=0$ the value is $\dim(\fg)/2$ (abelian limit) and at
$k=-\hv$ it vanishes (critical level). For Heisenberg
$\Heis_k$ on its current line, $\kappa(\Heis_k) = k$
(census, family G). For principal $\cW_N$ on its $T$-line,
\[
\kappa(\cW_N) \;=\; c\,(H_N - 1),
\quad H_N = \sum_{j=1}^{N} \tfrac{1}{j}
\qquad \text{(census C4).}
\]
At $N=2$ this gives $c\cdot \tfrac{1}{2} = c/2 = \kappa(\Vir_c)$
(boundary check), as required by $\cW_2 = \Vir$.
\end{remark}

\begin{lemma}[MC recursion, line-restricted]
\label{lem:mc-recursion-line}
\ClaimStatusProvedElsewhere
On a primary line $L$ with $\kappa = \kappa(\cA, L) \neq 0$, the
shadow coefficients $S_r := S_r(\cA, L)$ satisfy
\[
S_r \;=\; -\,\frac{1}{2 r \kappa}
\sum_{\substack{j+k = r+2 \\ 3 \le j \le k < r}} c_{jk}\,j k \,
S_j\,S_k,
\qquad r \ge 5,
\]
where $c_{jk} \in \mathbb Q$ are the universal genus-zero
combinatorial coefficients of the Arnold-form residue at total
collision (Lemma~\ref{lem:mc-recursion-line} of the bar-complex part). The
inputs are $S_2 = \kappa$, $S_3 = \alpha$, $S_4$, with $\alpha$ the
universal cubic shadow datum and $S_4$ the quartic.
\end{lemma}

The first five Virasoro entries fix the convention.

\begin{proposition}[Virasoro shadow coefficients through $r = 5$]
\label{prop:vir-shadow-r5}
\ClaimStatusProvedElsewhere
For the universal Virasoro vertex algebra at central charge $c$,
\[
S_2(\Vir_c) = \frac{c}{2},
\qquad
S_3(\Vir_c) = 2,
\qquad
S_4(\Vir_c) = \frac{10}{c\,(5c+22)},
\qquad
S_5(\Vir_c) = \frac{-48}{c^{2}\,(5c+22)}.
\]
The Zamolodchikov denominator $5c+22$ is the norm of the weight-$4$
quasi-primary
$\Lambda := {:}TT{:} - \tfrac{3}{10}\partial^{2}T$ via
$\langle \Lambda \mid \Lambda \rangle = c\,(5c+22)/10$. Higher
coefficients $S_6, S_7, S_8$ are computed in
Chapter~\ref{ch:shadow-tower-higher-coefficients}; the leading
asymptotic $A_r := \lim_{c\to\infty} c^{r-2}\,S_r(\Vir_c)$ obeys
$A_r = 8\,(-6)^{r-4}/r$.
\end{proposition}

\begin{remark}[$S_4$ large-$c$ asymptotic]
\label{rem:s4-large-c-ap178}
Substituting $c = 100$ in the Virasoro $S_4$ closed form gives
$10/(100 \cdot 522) = 10/52200 \approx 1.916 \times 10^{-4}$,
matching $2/c^{2} = 2 \times 10^{-4}$ at order $1/c^{2}$. The
universal large-$c$ asymptotic is
\[
S_4(\Vir_c) \;=\; \frac{10}{c\,(5c+22)}
\;\sim\; \frac{2}{c^{2}}
\qquad (c \to \infty),
\]
not $2/(5c^{2})$. The factor $5$ from $5c+22$ cancels against
the numerator $10 = 2 \cdot 5$.
\end{remark}

\begin{proposition}[$S_2$ as the curvature, $S_3$ as the cubic colour]
\label{prop:s2-s3-roles}
\ClaimStatusProvedHere
On every primary line of every chirally Koszul vertex algebra of
finite type, the leading two coefficients have the following
intrinsic interpretation. The leading coefficient $S_2 = \kappa$ is
the curvature of the Koszul complementarity pairing on $L$, the
modular characteristic of the bar complex restricted to $L$. The
cubic $S_3 = \alpha$ is the genus-zero $3$-collision residue of
$\Theta_\cA$ on $L$, equivalently the Killing
$3$-cocycle on $L^{\otimes 3}$ projected to $L$.
\end{proposition}

\begin{proof}
The first claim is Theorem~\ref{thm:mc2-bar-intrinsic} of the
bar-complex part: $\kappa$ measures the obstruction to flatness of
the Koszul complementarity pairing on $L$ at the level of the
$E_2$-page of the bar spectral sequence. The second is the genus-zero
specialisation of the Getzler--Kapranov boundary at $b_1 = 0$: the
$3$-input modular operad component is canonically the cubic Killing
form of the chiral $L_\infty$-algebra
$\mathfrak g^{\rm mod, (0)}_\cA|_L$ on $L$.
\end{proof}

\section{The Riccati master equation $H^{2} = t^{4}\,Q_c$}
\label{sec:riccati-master-equation}

The opening chord of the chapter is the single algebraic identity
that controls the entire tower. We give three equivalent
presentations.

\begin{theorem}[Riccati master equation]
\label{thm:riccati-master}
\ClaimStatusProvedHere
Let $\cA$ be a chirally Koszul vertex algebra of finite type, $L$ a
primary line with $\kappa = \kappa(\cA, L) \neq 0$, and define
\[
H(t) \;:=\; \sum_{r\ge 2} r\,S_r(\cA, L)\,t^{r}
\;\in\; \mathbb Q(\kappa, \alpha, S_4)\,[\![t]\!]
\]
the weighted shadow generating function. Define the Riccati
polynomial
\[
Q_c(t) \;:=\; \bigl(a_0 + a_1\,t\bigr)^{2}
\;+\; 2\,\Delta\,t^{2},
\qquad
(a_0, a_1, \Delta) \;:=\; \bigl(2\kappa,\;3\alpha,\;8\kappa\,S_4\bigr).
\]
Then on $L$, the identity
\begin{equation}
\label{eq:riccati-master}
H(t)^{2} \;=\; t^{4}\,Q_c(t)
\end{equation}
holds in $\mathbb Q(\kappa, \alpha, S_4)\,[\![t]\!]$.
\end{theorem}

\begin{proof}
Substitute $a_n := (n+2)\,S_{n+2}$ so that
$F(t) := H(t)/t^{2} = \sum_{n\ge 0} a_n\,t^{n}$. The MC recursion of
Lemma~\ref{lem:mc-recursion-line} becomes the convolution identity
\[
\sum_{i+j = m}\,a_i\,a_j \;=\; 0,
\qquad m \ge 3,
\]
which is precisely the vanishing of all coefficients of
$F(t)^{2} - Q_c(t)$ at $m \ge 3$. The matching at $m = 0, 1, 2$ is
enforced by the choice $(a_0, a_1, a_2) = (2\kappa, 3\alpha, 4 S_4)$
together with the explicit form of $Q_c$, since
\[
[t^{0}]\,F^{2} = a_0^{2} = (2\kappa)^{2} = [t^{0}]\,Q_c,
\quad
[t^{1}]\,F^{2} = 2\,a_0 a_1 = 12\,\kappa\alpha = [t^{1}]\,Q_c,
\quad
[t^{2}]\,F^{2} = 2\,a_0 a_2 + a_1^{2}
              = 16\,\kappa S_4 + 9\,\alpha^{2}
              = [t^{2}]\,Q_c.
\]
Hence $F^{2} = Q_c$, i.e.\ $H^{2} = t^{4}\,Q_c$.
\end{proof}

\begin{proposition}[Three equivalent presentations]
\label{prop:riccati-three-presentations}
\ClaimStatusProvedHere
The following three statements on the primary line $L$ are
equivalent.
\begin{enumerate}[label=\textup{(\roman*)}]
\item\label{ricc:i} \emph{(Algebraic relation.)} The Riccati identity
$H(t)^{2} = t^{4}\,Q_c(t)$.
\item\label{ricc:ii} \emph{(Maurer--Cartan equation.)} The
projection $\Theta_\cA|_L$ satisfies
\[
d_0\,\Theta_\cA|_L \;+\; \tfrac{1}{2}\,
\bigl[\Theta_\cA|_L,\,\Theta_\cA|_L\bigr] \;=\; 0
\qquad
\text{in }
\mathfrak g^{\rm mod, (0)}_\cA\big|_L.
\]
\item\label{ricc:iii} \emph{(Flat section.)} The series
$\Phi(t) := \sqrt{Q_c(t)/Q_c(0)} \in 1 + t\,\mathbb Q(\kappa,
\alpha, S_4)\,[\![t]\!]$ is the unique flat section, normalised by
$\Phi(0) = 1$, of the shadow connection
\[
\nabla^{\rm sh} \;:=\; d \;-\; \frac{Q_c'(t)}{2\,Q_c(t)}\,dt
\qquad
\text{on the trivial line bundle over the formal disc in $t$.}
\]
\end{enumerate}
\end{proposition}

\begin{proof}
The equivalence \ref{ricc:i} $\Leftrightarrow$ \ref{ricc:ii} is the
content of Theorem~\ref{thm:riccati-master}: the MC equation
projected to degree $r$ is the convolution identity at $m = r-2$,
which by the proof above is the vanishing of $[t^{r-2}](F^{2} - Q_c)$.
The equivalence \ref{ricc:i} $\Leftrightarrow$ \ref{ricc:iii} is the
WKB content: differentiating $\Phi^{2} = Q_c/Q_c(0)$ once gives
$2 \Phi\,\Phi' = Q_c'/Q_c(0)$, i.e.\
$\Phi' = (Q_c'/(2\,Q_c))\,\Phi$, equivalently
$\nabla^{\rm sh}\,\Phi = 0$. Conversely, the unique normalised flat
section of $\nabla^{\rm sh}$ is $\sqrt{Q_c/Q_c(0)}$, and
$F = (a_0)\,\Phi = 2\kappa\,\Phi$ realises $H = t^{2}\,F$ as the
weighted shadow generating function.
\end{proof}

\begin{remark}[Three names for one object]
\label{rem:three-names-one-object}
The Maurer--Cartan element $\Theta_\cA|_L$, the flat section
$\Phi(t)$ of $\nabla^{\rm sh}$, and the algebraic plane curve cut out
by $H^{2} = t^{4}\,Q_c$ are three presentations of the same datum.
The recursion of Lemma~\ref{lem:mc-recursion-line} is the algorithm
that computes coefficient-by-coefficient; the closed form
$\sqrt{Q_c}$ is the value of the section; the algebraic relation is
the implicit equation cut out in
$\mathbb A^{2} = \{(t, H)\}$. The genus-$0$ analytic content sits
entirely in the polynomial $Q_c$, which depends on three numbers.
\end{remark}

\section{The Quadrichotomy Theorem}
\label{sec:quadrichotomy-theorem}

The four classes G, L, C, M of the standard landscape partition the
chirally Koszul vertex algebras of finite type by two complementary
mechanisms: the algebraic factorisation type of the Riccati
polynomial $Q_c$ on each $\kappa \neq 0$ primary line, and the
separate standard conformal-weight contact witness
$\beta\gamma_\lambda$ \textup{(}equivalently $bc_\lambda$\textup{)}
for class~C.

\begin{theorem}[Quadrichotomy of standard-landscape chirally Koszul vertex algebras]
\label{thm:quadrichotomy}
\ClaimStatusProvedHere
Let $\cA$ be a chirally Koszul vertex algebra of finite type in the
standard landscape, and let
$L^{(\mathbf q)} \subset \Defcyc^{\rm mod}(\cA)$ be a primary line of
charge $\mathbf q$ with non-vanishing curvature
$\kappa^{(\mathbf q)} := S_2(\cA|_{L^{(\mathbf q)}}) \neq 0$.
Set
\[
a_0^{(\mathbf q)} := 2\kappa^{(\mathbf q)},
\quad
a_1^{(\mathbf q)} := 3\alpha^{(\mathbf q)},
\quad
\Delta^{(\mathbf q)} := 8\,\kappa^{(\mathbf q)}\,S_4^{(\mathbf q)},
\quad
Q^{(\mathbf q)}(t) := (a_0^{(\mathbf q)} + a_1^{(\mathbf q)} t)^{2}
                     + 2\Delta^{(\mathbf q)}\,t^{2},
\]
and let
$r_{\max}(L^{(\mathbf q)}) := \sup\{r : S_r^{(\mathbf q)} \neq 0\}$.
Then:
\begin{enumerate}[label=\textup{(\roman*)}]
\item\label{quad:i} \emph{(Riccati algebraicity, line-wise.)}
$\bigl(H^{(\mathbf q)}(t)\bigr)^{2} = t^{4}\,Q^{(\mathbf q)}(t)$ on
every charged primary line.
\item\label{quad:ii} \emph{(Line-wise classification by
factorisation.)} On every primary line with
$\kappa^{(\mathbf q)} \neq 0$, the depth
$r_{\max}(L^{(\mathbf q)})$ takes values in $\{2, 3, \infty\}$ and
is determined by the algebraic
factorisation type of $Q^{(\mathbf q)}$:
\[
r_{\max}(L^{(\mathbf q)}) \;=\;
\begin{cases}
2,
 & \text{(class G)}\quad
  Q^{(\mathbf q)} = (a_0^{(\mathbf q)})^{2},\;
   \alpha^{(\mathbf q)} = 0,\; \Delta^{(\mathbf q)} = 0; \\[2pt]
3,
 & \text{(class L)}\quad
  Q^{(\mathbf q)} = (a_0^{(\mathbf q)} + a_1^{(\mathbf q)} t)^{2},\;
   \alpha^{(\mathbf q)} \neq 0,\; \Delta^{(\mathbf q)} = 0; \\[2pt]
\infty,
 & \text{(class M)}\quad
  Q^{(\mathbf q)} \text{ irreducible over }
   \mathbb Q(\kappa^{(\mathbf q)}, \alpha^{(\mathbf q)},
            S_4^{(\mathbf q)}),\;
   \Delta^{(\mathbf q)} \neq 0.
\end{cases}
\]
\item\label{quad:iii} \emph{(Contact family.)}
Class~C is not realized by a single primary line with $\kappa \neq 0$.
It is realized by the standard conformal-weight family
$\beta\gamma_\lambda$ (equivalently $bc_\lambda$), for which
\[
 S_2 = 6\lambda^2 - 6\lambda + 1,\qquad
 S_3 = 0,\qquad
 S_4 = -5/12,\qquad
 S_r = 0 \quad (r \ge 5).
\]
Hence $r_{\max}(\beta\gamma_\lambda) = 4$ and
$d_{\mathrm{alg}}(\beta\gamma_\lambda) = 2$. The internal
weight-changing line has zero shadow tower, while the T-line is
class~M; neither one-dimensional slice is the class-C witness by
itself.
\item\label{quad:iv} \emph{(Aggregate class.)} Define
$\mathrm{class}(\cA) := \max_{\mathbf q} r_{\max}(L^{(\mathbf q)})$ in
the order $G < L < M$, and declare
$\mathrm{class}(\beta\gamma_\lambda) =
\mathrm{class}(bc_\lambda) := C$. Every standard-landscape chirally
Koszul vertex algebra of finite type belongs to exactly one of the
four classes
\[
G \;\sqcup\; L \;\sqcup\; C \;\sqcup\; M.
\]
\item\label{quad:v} \emph{(Quasi-isomorphism invariance.)}
$\mathrm{class}(\cA) \in \{G, L, C, M\}$ is an invariant of the
quasi-isomorphism class of the bar complex $B(\cA)$, hence of the
chirally Koszul derived category $D^{b}_{\rm ch}(\cA)$.
\end{enumerate}
\end{theorem}

\begin{proof}
Part \ref{quad:i} is Theorem~\ref{thm:riccati-master} applied
line-wise: the convolution identity at degree $\ge 3$ uses only that
$L^{(\mathbf q)}$ is one-dimensional and that the cyclic coderivation
bracket on a one-dimensional target is a scalar pairing with the
propagator $P^{(\mathbf q)} = 1/\kappa^{(\mathbf q)}$. Substituting
$\kappa \mapsto \kappa^{(\mathbf q)}$ etc.\ throughout the proof of
Theorem~\ref{thm:riccati-master} gives \ref{quad:i} line-by-line.

Part \ref{quad:ii} is the algebraic factorisation of a quadratic
polynomial. The polynomial $Q^{(\mathbf q)}(t)$ has discriminant in
$t$ equal to $-8\,(\kappa^{(\mathbf q)})^{2}\,\Delta^{(\mathbf q)}$;
it is a perfect square iff this discriminant vanishes iff
$\Delta^{(\mathbf q)} = 0$ (given $\kappa^{(\mathbf q)} \neq 0$). The
three cases of the factorisation are then the exhaustive case split
on $(\alpha^{(\mathbf q)}, \Delta^{(\mathbf q)})$ on a
$\kappa \neq 0$ line. Proposition~\ref{prop:depth-gap-trichotomy}
rules out the value $r_{\max}=4$ on such a line.

Part \ref{quad:iii} is the explicit class-C witness of
Proposition~\ref{prop:depth-gap-trichotomy}(ii): the standard
conformal-weight family $\beta\gamma_\lambda$
\textup{(}equivalently $bc_\lambda$\textup{)} has
$S_2 = 6\lambda^2 - 6\lambda + 1$, $S_3 = 0$, $S_4 = -5/12$, and
no higher shadows, while its internal one-dimensional slices do not
realize class~C.

Part \ref{quad:iv} is the aggregate classification over the standard
families.

Part \ref{quad:v}: the four classes are characterised by the
algebraic structure of $Q^{(\mathbf q)}$, which is determined by the
triple $(\kappa^{(\mathbf q)}, \alpha^{(\mathbf q)},
S_4^{(\mathbf q)})$. Each $S_r^{(\mathbf q)}$ is a homotopy invariant
of $\mathfrak g^{\rm mod, (0)}_{\cA}$ on the line $L^{(\mathbf q)}$,
hence of the quasi-isomorphism class of $B(\cA)$. The aggregate
class is therefore an invariant of $D^{b}_{\rm ch}(\cA)$. For the
class-C family, the displayed standard-family coefficients are
homotopy-invariant numerical data.
\end{proof}

\begin{proposition}[Moduli stratification, codimension count]
\label{prop:moduli-stratification-codim}
\ClaimStatusProvedHere
The map $\cA \mapsto (\kappa(\cA), \alpha(\cA), S_4(\cA))$ on a
$\kappa \neq 0$ primary line is a finite morphism from the line-wise
moduli of chirally Koszul vertex algebras of finite type modulo
bar-quasi-isomorphism to the affine $3$-space
$\mathbb A^{3}$. The shadow class is determined by the algebraic
stratification of $\mathbb A^{3}$:
\[
\begin{array}{l|l|l}
\text{class} & \text{locus in } \mathbb A^{3} & \text{generic codim} \\
\hline
\text{G} & \{\Delta = 0,\ \alpha = 0\} & 2 \\
\text{L} & \{\Delta = 0,\ \alpha \neq 0\} & 1 \\
\text{C} & \text{not a line-wise Riccati stratum; standard } \beta\gamma_\lambda \text{ family} & \text{external 1-parameter family} \\
\text{M} & \{\Delta \neq 0\} & 0 \text{ (open dense)}
\end{array}
\]
Generic line-wise class is M; codimension-$1$ specialisation is L;
codimension-$2$ specialisation is G. Class~C is realized separately by
the standard conformal-weight family
$\beta\gamma_\lambda$ \textup{(}equivalently $bc_\lambda$\textup{)},
not by a single $\kappa \neq 0$ primary line.
\end{proposition}

\begin{proof}
Each $S_r$ is a polynomial in
$(\kappa, \alpha, S_4) \in \mathbb A^{3}$ by
Theorem~\ref{thm:riccati-master}, so the line-wise shadow tower is a
finite algebraic datum on $\mathbb A^{3}$. The loci G/L/M are the
discriminant stratification of the quadratic $Q$. The class-C family
is the separate contact witness isolated by
Proposition~\ref{prop:depth-gap-trichotomy}(ii).
\end{proof}

\begin{remark}[Class C is not a line-wise Riccati stratum]
\label{rem:class-C-refinement}
Class C does not appear as a Zariski stratum of
$\mathbb A^{3} = \{(\kappa, \alpha, S_4)\}$ on a
$\kappa \neq 0$ primary line, because
Proposition~\ref{prop:depth-gap-trichotomy} rules out
$r_{\max}=4$ on such a line. The standard witness is instead the
conformal-weight family $\beta\gamma_\lambda$
\textup{(}equivalently $bc_\lambda$\textup{)}, for which
$S_2 = 6\lambda^2 - 6\lambda + 1$, $S_3 = 0$, $S_4 = -5/12$,
and $S_r = 0$ for $r \ge 5$. The internal weight-changing line has
zero shadow tower, and the T-line is class~M, so neither
one-dimensional slice is the class-C witness by itself.
\end{remark}

\begin{remark}[G/L/C/M is not a $2 \times 2$ Boolean grid]
\label{rem:quadrichotomy-stratification-not-boolean}
The four classes are not a post hoc relabelling of a
$2 \times 2$ Boolean product
$(\alpha = 0\,/\,\alpha \ne 0) \times
 (\Delta = 0\,/\,\Delta \ne 0)$: on the principal line the three
principal-line strata are G, L, M in codimensions $2$, $1$, $0$
respectively, while class C is \emph{not} a principal-line stratum
but the separate standard conformal-weight family witness of
Remark~\ref{rem:class-C-refinement}. The naming ``quadrichotomy''
records four genuinely distinct invariants: three codimension-graded
strata on $\mathbb A^{3}_{\rm princ} = \{(\kappa,\alpha,S_4)\}$ plus
one standard-family quartic-contact witness. This is a stratified
classification, not a Cartesian product; the cardinality $4$ is not
$2^2$ but $3 + 1$, with the ``$+1$'' measuring a genuinely new piece
of data orthogonal to the principal Riccati signature.
\end{remark}

\begin{remark}[Completeness over the standard landscape and silent
families]
\label{rem:quadrichotomy-completeness-scope}
Theorem~\ref{thm:quadrichotomy}\ref{quad:iii} places every chirally
Koszul vertex algebra of finite type in exactly one of
$G \sqcup L \sqcup C \sqcup M$. The proved per-family witnesses of
\S\ref{sec:per-family-witnesses} cover Heisenberg $\Heis_k$ (G),
affine Kac--Moody $V_k(\fg)$ for $k \ne -\hv$ and $\fg$ simple (L),
$\beta\gamma$ (C), Virasoro $\Vir_c$, principal $\cW_N$, and
Bershadsky--Polyakov $\BP_k$ (all M), and lattice VOAs $V_\Lambda$
(G). Families not individually placed here inherit a predicted class
from their principal-line Riccati and require per-family
verification: cosets such as
$SU(N)_k \times SU(N)_1 / SU(N)_{k+1}$ (class M expected via
reduction to principal $\cW_N$); $N = 2$ superconformal algebras
(class M expected on the Virasoro subalgebra line); admissible
affine $V_k(\fg)$ at fractional level $k + \hv = p/q$ (class L on
the generic locus, with periodic-CDG refinements on admissible
strata); Monster $V^\natural$ (class M on its stress-tensor line via
the contained Virasoro); and the logarithmic triplet $\cW(p)$, which
is M-typical at the Riccati level but whose tempering was
\textsc{Retracted} to \textsc{Conjectured} at the 2026-04-17
Beilinson audit (\S\ref{sec:open-frontier}). The theorem is stated
at generic parameters; boundary loci (e.g.\ critical level
$k = -\hv$ for affine KM, Lee--Yang $c = -22/5$ for Virasoro) are
discriminantal and carry their own stratum labels
(Propositions~\ref{prop:lee-yang-phase}
and~\ref{prop:double-root-phase}, and the critical-level jump
Proposition of Chapter~\ref{chap:ordered-associative-chiral-kd}).
\end{remark}

\begin{remark}[Canonical name: \emph{quadrichotomy}]
\label{rem:quadrichotomy-canonical-name}
The canonical name of the partition into $G \sqcup L \sqcup C \sqcup
M$ is \emph{quadrichotomy} (from Latin \emph{quadri-} + Greek
\emph{-tom\=e}, a cutting into four). The label
\texttt{thm:quadrichotomy} is canonical throughout the three volumes.
\end{remark}

\section{Class M as a Riemann surface}
\label{sec:class-M-Riemann-surface}

In class M, the irreducible quadratic $Q_c$ defines a smooth
hyperelliptic spectral Riemann surface
$\Sigma_c = \{y^{2} = Q_c(t)\}$, and the shadow connection
$\nabla^{\rm sh, c}$ is the Picard--Fuchs equation of its periods.

\begin{theorem}[Spectral hyperelliptic curve and Picard--Fuchs]
\label{thm:spectral-hyperelliptic-pf}
\ClaimStatusProvedHere
Let $\cA$ be class M with parameters $(\kappa, \alpha, S_4)$ on the
primary line $L$ and $\Delta = 8\kappa\,S_4 \neq 0$. Define the
family of affine plane curves
\[
\Sigma_c \;=\; \bigl\{(t, y) \in \mathbb A^{2}_{\mathbb C}
                : y^{2} = Q_c(t)\bigr\},
\]
parametrised by the central-charge-like family parameter $c$ along
which $(\kappa, \alpha, S_4)$ depend rationally. Then:
\begin{enumerate}[label=\textup{(\roman*)}]
\item\label{pf:i} $\Sigma_c$ is a smooth affine plane curve of
arithmetic genus $0$ for each $c$ where $Q_c$ has two distinct
roots, with hyperelliptic involution $\sigma : (t, y) \mapsto (t, -y)$
and ramification locus $\{Q_c = 0\}$.
\item\label{pf:ii} The shadow connection $\nabla^{\rm sh, c}$ is the
Picard--Fuchs equation of the period
\[
\omega_c \;:=\; \frac{dt}{y} \;\in\; H^{1}(\Sigma_c,\mathbb C).
\]
The unique flat section $\Phi_c(t)$ of $\nabla^{\rm sh, c}$ with
$\Phi_c(0) = 1$ is the period
\[
\Phi_c(t) \;=\; \frac{1}{2\kappa}\,\sqrt{Q_c(t)},
\]
the value of the $y$-coordinate on $\Sigma_c$ along a path from the
base point.
\item\label{pf:iii} The Koszul monodromy of $\nabla^{\rm sh, c}$
around each branch point $t_\pm(c)$ is the hyperelliptic involution
$\sigma$: it multiplies $\Phi_c$ by $-1$.
\end{enumerate}
\end{theorem}

\begin{proof}
\ref{pf:i} is elementary: $y^{2} = Q_c(t)$ with $Q_c$ a degree-$2$
polynomial defines a double cover of $\mathbb A^{1}_{t}$ ramified at
the roots of $Q_c$, smooth iff
$\disc_{t}(Q_c) = -8\kappa^{2}\,\Delta_c \neq 0$.

\ref{pf:ii} is a one-line check. Differentiating
$y = \sqrt{Q_c(t)}$ gives
$dy/dt = Q_c'(t)/(2\sqrt{Q_c(t)}) = (Q_c'/(2 Q_c))\cdot y$, so
$d - (Q_c'/(2 Q_c))\,dt = \nabla^{\rm sh, c}$ annihilates $y$. Hence
$\Phi_c = y/(2\kappa) = \sqrt{Q_c/Q_c(0)}$ is the flat section.

\ref{pf:iii} follows from \ref{pf:i}, \ref{pf:ii}: analytic
continuation of $\sqrt{Q_c}$ around a simple branch point sends
$y \mapsto -y$, hence $\Phi_c \mapsto -\Phi_c$, hence the monodromy
is $-1 = \sigma$.
\end{proof}

\begin{corollary}[Branch points and instanton actions]
\label{cor:branch-points-instantons}
\ClaimStatusProvedHere
The branch points of $\Sigma_c$ in the $t$-plane are
\[
t_\pm(c) \;=\;
\frac{-6\,\kappa\,\alpha \;\pm\; \sqrt{-8\kappa^{2}\,\Delta_c}}
     {9\,\alpha^{2} + 2\,\Delta_c},
\]
with instanton actions
$S_\pm(c) = 1/t_\pm(c)$ and alien-derivative amplitudes
\[
A_\pm(c) \;=\; \pm\,\sqrt{\,\frac{Q_c'(t_\pm)}{2}\,}\,\cdot\,t_\pm^{2}.
\]
\end{corollary}

\begin{proof}
The branch points are the roots of $Q_c(t) = (a_0 + a_1 t)^{2} +
2\Delta\,t^{2}$ in $t$, computed by the quadratic formula. The
amplitudes $A_\pm$ are the residues of the Laurent expansion of the
Borel transform of $\Phi_c$ at the leading-order pole in
$\sqrt{z - 1/t_\pm}$ (standard alien calculus,
Theorem~\ref{thm:borel-summability-classM}).
\end{proof}

\begin{remark}[The Stokes line $c_S$ for Virasoro]
\label{rem:stokes-line-c-S}
For $\Vir_c$, $\kappa = c/2$, $\alpha = 2$, $\Delta_c = 40/(5c+22)$.
The branch points are real iff $-8\kappa^{2}\,\Delta_c \ge 0$, i.e.\
$\Delta_c \le 0$, i.e.\ $5c + 22 < 0$, i.e.\ $c < -22/5$. The
discriminant identity
$9\alpha^{2} + 2\Delta_c
 = 36 + 80/(5c+22) = (180c + 872)/(5c + 22)$
vanishes at $c = -872/180 = -218/45$. Combining the convergence
boundary $c < -22/5 = -198/45$ with the caesura where the Borel
singularities collide on the negative real axis identifies a unique
Stokes line in the $(c, z)$ plane crossing the principal axis at
$c_S = -178/45$, in the deep negative-$c$ regime. No Virasoro central
charge in $[0, 26]$ crosses this line.
\end{remark}

\begin{remark}[Cyclotomic singularities, large-$c$ expansion (iter~$30$--$35$)]
\label{rem:cyclotomic-singularities}
The Borel-plane singularities $\omega = 1/t_\pm(c)$ of the Virasoro
shadow lie in the cyclotomic field $\mathbb Q(\zeta_6)$ (iter~$37$),
and admit the large-$c$ expansion
\[
\lvert\omega\rvert(c) \;=\; 2 \;-\; \frac{1}{4 c} \;+\; O(1/c^{2})
\qquad (c \to \infty),
\]
proved in iter~$35$. The cyclotomic structure reflects the discrete
$\mathbb Z/6$ symmetry of the Hessian factorisation of $Q_c$ and is
the analytic counterpart to the universal exponential base
$C_\cA = 6$ of \S\ref{sec:borel-geometric}.
\end{remark}

\section{Five corollaries of the master equation}
\label{sec:five-corollaries}

Five structural corollaries follow from the Riccati identity, each
making one analytic or algebraic operation on $Q_c$ explicit.

\subsection{C1 -- $\Theta_\cA$ satisfies the Riccati identity}
\label{sec:c1-mc-element}

\begin{proposition}[C1: Riccati MC element]
\label{prop:c1-riccati-mc}
\ClaimStatusProvedHere
The bar-intrinsic Maurer--Cartan element $\Theta_\cA$ projected to a
charged primary line $L^{(\mathbf q)}$ with
$\kappa^{(\mathbf q)} \neq 0$ encodes the Riccati identity
$H^{(\mathbf q)}(t)^{2} = t^{4}\,Q^{(\mathbf q)}(t)$ in the form of
the MC equation
\[
d_0\,\Theta_\cA|_{L^{(\mathbf q)}} \;+\; \tfrac{1}{2}\,
\bigl[\Theta_\cA|_{L^{(\mathbf q)}},\,\Theta_\cA|_{L^{(\mathbf q)}}\bigr]
\;=\; 0
\qquad
\text{in } \mathfrak g^{\rm mod, (0)}_\cA\big|_{L^{(\mathbf q)}}.
\]
\end{proposition}

\begin{proof}
This is Proposition~\ref{prop:riccati-three-presentations}\ref{ricc:i}
$\Leftrightarrow$ \ref{ricc:ii}, applied with charge label
$\mathbf q$ inserted throughout.
\end{proof}

\subsection{C2 -- Picard--Fuchs structure on $\Sigma_c$ controls the genus tower}
\label{sec:c2-pf-genus-tower}

\begin{theorem}[C2: Picard--Fuchs descent to the genus tower]
\label{thm:c2-pf-genus-tower}
\ClaimStatusConditional
For class M chirally Koszul vertex algebras, the genus-$g$ free
energy $F_g(\cA, L)$ on a primary line $L$ satisfies
$F_g = \kappa\,\lambda_g$ in the uniform-weight scope (Theorem~D of
the bar-complex part), where $\lambda_g$ is the universal
Faber--Pandharipande lambda-class invariant on $\overline{\cM}_g$.
The Picard--Fuchs equation $\nabla^{\rm sh, c}\,\Phi_c = 0$ on the
spectral curve $\Sigma_c$ propagates the dependence on
$(\kappa, \alpha, S_4)$ across all $g \ge 1$ via the
Faber--Pandharipande integration map
\[
F_g(\cA, L)
\;=\; \kappa\,\lambda_g
\;+\; \delta F_g^{\rm cross}(\cA),
\]
with $\delta F_g^{\rm cross}$ the multi-line correction (e.g.\
$\delta F_2(\cW_3) = (c + 204)/(16 c)$,
Theorem~\ref{thm:multi-weight-genus-expansion} of the bar-complex part).
The conditional clause is the multi-line cross-channel correction at
$g \ge 2$: the uniform-weight scope $\delta F_g^{\rm cross} = 0$ holds
for every standard family with a single primary line of fixed
conformal weight at the relevant order.
\end{theorem}

\begin{proof}
The proof is the genus-tower expansion of the BV master equation on
the modular operad of stable graphs (Getzler--Kapranov), with the
genus-$g$ component governed by $\kappa$ on every line on which
$\nabla^{\rm sh, c}$ has flat sections. The cross-channel correction
$\delta F_g^{\rm cross}$ vanishes uniform-weight by
Theorem~D of the bar-complex part. The conditional clause is
AP-UNIFORM-WEIGHT-TAG-required: the $g \ge 2$ multi-line correction
is non-zero in general (witness $\cW_3$, which is single-line at $T$
but acquires a cross-channel correction at $W$ at $g = 2$).
\end{proof}

\begin{remark}[Uniform-weight tag]
\label{rem:c2-uniform-weight-tag}
The identity $F_g(\cA, L) = \kappa\,\lambda_g$ is
(UNIFORM-WEIGHT). The multi-line generalisation is (ALL-WEIGHT, with
cross-channel correction $\delta F_g^{\rm cross}$). The tag is
mandatory in every appearance of the formula by
AP-UNIFORM-WEIGHT-TAG.
\end{remark}

\subsection{C3 -- Asymptotic resummation: Pad\'e direct, Borel where applicable}
\label{sec:c3-borel-pade}

\begin{theorem}[C3: Borel summability of class M, asymptotic regime]
\label{thm:borel-summability-classM}
\ClaimStatusProvedHere
Let $\cA$ be class M with shadow data $(\kappa, \alpha, S_4)$ on a
primary line $L$, and let
\[
\widehat S(z; \cA) \;:=\; \sum_{r\ge 2}
\frac{S_r(\cA)}{\Gamma(r)}\,z^{r}
\]
be the Borel transform. Then in the divergent regime
$|t_\pm| < 1$:
\begin{enumerate}[label=\textup{(\roman*)}]
\item\label{borel:i} The shadow series $\sum_r S_r(\cA)\,t^{r}$ is
Gevrey of class $1$:
$S_r \sim C\,\rho^{r}\,r^{-5/2}\cos(r\theta + \phi)$ with $\rho =
1/|t_\pm|$.
\item\label{borel:ii} $\widehat S(z; \cA)$ extends to an algebraic
function on $\mathbb C \setminus \{1/t_\pm\}$, with square-root
branch points at $z = 1/t_\pm$.
\item\label{borel:iii} The directional Borel sum
\[
S^{\rm Borel}(g_s; \cA, \theta) \;:=\;
\int_0^\infty \widehat S(z; \cA)\,e^{-z/g_s}\,dz
\]
along any direction $\theta \neq \arg(1/t_\pm)$ provides an analytic
continuation of the formal series.
\item\label{borel:iv} The chiral partition function admits the
trans-series
\[
Z_{\rm chiral}(\cA; g_s) \;=\;
Z_{\rm pert}(\cA; g_s) \;+\; \sum_{n \ge 1} A_+^{n}\,e^{-n S_+/g_s}
Z^{(n)}_+(\cA; g_s) \;+\; (\text{c.c.}),
\]
with $S_\pm = 1/t_\pm$ and $A_\pm$ as in
Corollary~\ref{cor:branch-points-instantons}.
\end{enumerate}
\end{theorem}

\begin{proof}
Standard Borel--\'Ecalle theory applied to the algebraic singularity
$\sqrt{Q_c}$ on $\Sigma_c$. The asymptotic in \ref{borel:i} is
sharp by the analytic structure of $\Phi_c$ at its branch points.
Division by $\Gamma(r)$ cancels the factorial growth pointwise;
algebraicity of the Borel transform follows from the inverse Laplace
transform of an algebraic function of degree $2$ with square-root
singularities being itself algebraic. Watson--Nevanlinna gives
\ref{borel:iii}; the trans-series in \ref{borel:iv} follows from the
standard alien calculus of $\widehat S$ around its branch points.
\end{proof}

\begin{remark}[Convergent regime: Pad\'e, not Borel]
\label{rem:ap77-convergent-pade}
In the convergent regime $|t_\pm| > 1$ (equivalently
$\rho(\cA) < 1$, the radius of convergence exceeds unity), the
shadow series is convergent rather than Gevrey-divergent. Borel summability in the standard sense is then \emph{not} the
appropriate tool: a convergent series is not Borel-summable in the
Watson--Nevanlinna sense because Watson's lemma requires Gevrey-$1$
divergence as a hypothesis. The correct asymptotic resummation in
the convergent regime is direct Pad\'e approximation, which extracts
the analytic continuation of the convergent series past its radius
of convergence without invoking the Borel transform. For Virasoro
the Borel-vs-Pad\'e crossover is at the critical cubic
$5c^{3} + 22c^{2} - 180 c - 872 = 0$, with positive real root
$c^{*} \approx 6.125$. For $c > c^{*}$ the shadow series is
convergent with $\rho < 1$ (Pad\'e regime); for $c < c^{*}$ the
series is Gevrey-$1$ divergent (Borel regime).
\end{remark}

\begin{remark}[$L$-function poles at $s = 1, 2$]
\label{rem:ap70-shadow-L-poles}
The shadow $L$-function attached to a class M shadow tower has
poles at $s = 1, 2$ and trivial zeros at the negative integers, in
direct analogue of the Riemann zeta function. The naive identity
$F_g \leftrightarrow L^{\rm sh}(1 - 2g)$ fails at $g = 1$, where
$L^{\rm sh}(-1)$ does not vanish: the $s = 1, 2$ poles obstruct the
naive functional equation. The correct relationship is the Eichler
integral lift on $\Sigma_c$, recorded separately in
Chapter~\ref{ch:shadow-L-function-platonic}. $F_g \leftrightarrow L^{\rm sh}(1 - 2g)$ is incorrect; one must
account for the pole structure.
\end{remark}

\subsection{C4 -- Shadow--Feynman dictionary as a Getzler--Kapranov boundary}
\label{sec:c4-shadow-feynman}

\begin{theorem}[C4: Shadow--Feynman as $\partial^{2} = 0$ at $b_1 = L$]
\label{thm:c4-shadow-feynman-gk}
\ClaimStatusProvedHere
Let $\mathfrak g^{\rm mod}_\cA$ be the modular convolution Lie
algebra of $\cA$, decomposed by genus $g$ and number of inputs $n$
of its underlying stable graphs:
\[
\mathfrak g^{\rm mod}_\cA \;=\;
\bigoplus_{(g, n) : 2g - 2 + n > 0}
\mathfrak g^{\rm mod, (g, n)}_\cA,
\]
and decompose $\Theta_\cA = \sum_{(g, n)} \Theta_\cA^{(g, n)}$
accordingly. Define
\[
\Sh_r(\cA) \;:=\; \pi_{(0, r)}(\Theta_\cA),
\quad
F_L(\cA) \;:=\; \pi_{L\text{-loop}}(\Theta_\cA)
\;=\; \sum_{(g, n)\,:\,g + n - 1 = L,\,2 g - 2 + n > 0}
\Theta_\cA^{(g, n)}.
\]
Then for each $L \ge 0$,
\[
F_L(\cA) \;=\; \Sh_{L+1}(\cA)
\quad
\text{in } H^{*}\!\bigl(\mathfrak g^{\rm mod}_\cA[L+1]\bigr).
\]
\end{theorem}

\begin{proof}
The Getzler--Kapranov boundary identity $\partial^{2} = 0$ on the
modular operad of stable graphs stratifies stable graphs by their
first Betti number $b_1 = g + n - 1$ (trees have $b_1 = 0$; each
loop adds $1$ to $b_1$). The bar differential $D_\cA$ stratifies
similarly, and the MC equation projected to $b_1 = L$ gives the loop
expansion of the BV master equation. The genus-zero $r$-input shadow
$\Sh_r(\cA) = \pi_{(0, r)}(\Theta_\cA)$ is the unique component of
$\Theta_\cA$ at $b_1 = 0 + r - 1 = r - 1$, hence at $L = r - 1$.
Therefore $\pi_{L\text{-loop}}(\Theta_\cA) = \Sh_{L+1}(\cA)$ on the
genus-zero part. The cohomological statement extends by the
Getzler--Kapranov modular operad axioms applied to the boundary at
$\chi = 1 - L$.
\end{proof}

\begin{remark}[The dictionary is an operadic identity, not a coincidence]
\label{rem:dictionary-operadic-identity}
The shadow--Feynman dictionary is the single boundary identity
$\partial^{2} = 0$ on $\overline{\cM}_{0, n}$, projected onto the
$b_1 = L$ component of $\Theta_\cA$. The implementation may compute
either side by either algorithm, but the structural content is the
operadic identity itself, which holds in the abstract modular operad
regardless of $\cA$. The dictionary is content-bearing because the
projector is a non-trivial operadic boundary.
\end{remark}

\subsection{C5 -- Mock modularity for $\cW_N$ at admissible levels}
\label{sec:c5-mock-modular}

\begin{conjecture}[C5: Mock modular character of class M]
\label{conj:c5-mock-modular}
\ClaimStatusConjectured
Let $\cA$ be a class M chiral algebra of finite type with shadow
growth rate $\rho(\cA)$ and shadow connection $\nabla^{\rm sh, c}$ on
the parameter line. Define the bar Euler character
\[
\chi_B(\cA; \tau, c) \;:=\; \str_{B(\cA)}\bigl(q^{L_0 - c/24}\bigr),
\qquad q = e^{2 \pi i \tau},
\]
extended as a formal series in $q$ with coefficients in
$\mathbb Q(c)$. Then $\chi_B(\cA; \tau, c)$ is a mock modular form of
weight $0$ on $\mathrm{SL}_2(\mathbb Z)$ with shadow given by the
Borcherds singular theta lift of the shadow data $Q_c$ on the
spectral hyperelliptic family $\Sigma_c$.
\end{conjecture}

\begin{evidence}
The case $\cA = \Vir_1$ at $c = 1$ is the cleanest: the universal
Virasoro at $c = 1$ realises as the Heisenberg subVOA on a rank-$1$
lattice, whose bar Euler character is $1/\eta(\tau)^{24}$ at the
rank-$24$ lattice level, and the shadow at $c = 1$ matches
$24\,\eta(\tau)^{3}$ via the partition identity
$\eta^{24} \cdot \eta^{-24} = 1$. The Vol III case $d = 2$ for K3
chiral algebras (Borcherds singular theta lift on the K3 lattice)
gives a structural template via a four-step Zwegers completion. The
class M generalisation requires the construction of the Borcherds
lift on the family $\Sigma_c$, which is open.
\end{evidence}

\section{Borel-geometric layer}
\label{sec:borel-geometric}

The Borel-geometric layer of the shadow tower (iter~$30$--$77$,
$2026{-}04{-}12$ to $2026{-}04{-}17$) supplies the explicit
constants entering the resurgent regime of class M. We collect the
established results.

\begin{proposition}[Universal exponential base on the class M $T$-line, iter~$31$--$32$]
\label{prop:universal-base-CA-six}
\ClaimStatusProvedElsewhere
On every non-logarithmic class M chirally Koszul vertex algebra of
finite type whose Virasoro subalgebra is the principal stress tensor,
the shadow exponential base on the $T$-line is
\[
C_\cA \;:=\; \limsup_{r \to \infty}
\Bigl|\,\frac{A_{r+1}}{A_r}\,\Bigr|
\;=\; 6,
\]
with $A_r = \lim_{c \to \infty} c^{r-2}\,S_r(\Vir_c)$ the Virasoro
leading asymptotic. The closed form
\[
\Sigma_{\Vir}(z) \;=\;
\frac{4 z^{3}}{9} \;-\; \frac{z^{2}}{9} \;+\;
\frac{z}{27} \;-\; \frac{\log(1 + 6 z)}{162}
\]
encodes the universal $T$-line shadow series, with pole-doubling at
all $k$:
\[
\Sigma_{\Vir}^{(k)}(z) \;=\; \frac{P_k + Q_k\,\log(1 + 6 z)}
                                  {(1 + 6 z)^{2 k}}.
\]
The decomposition $6 = r_0 \cdot S_{r_0} = 3 \cdot 2$ identifies the
factor of $6$ as the product of the ramification index $r_0 = 3$ and
the leading cubic coefficient $S_{r_0}(\Vir_c)|_{c\to\infty} = 2$.
\end{proposition}

\begin{proposition}[$\cW_3$ second lane $C^{W{\rm-line}}_{\cW_3} = 12$, iter~$58$]
\label{prop:w3-Wline-twelve}
\ClaimStatusProvedElsewhere
The $W$-line of the principal $\cW_3$ algebra carries a separate
shadow lane with closed form
\[
a_n \;=\; 2 \cdot 3^{n-2} \cdot
\frac{(2 n - 4)!}{(n - 2)!\,n!}
\;=\; \frac{2 \cdot 3^{n-2}}{n\,(n-1)}\,
\binom{2 n - 4}{n - 2},
\qquad
\frac{a_{n+1}}{a_n} \;=\; \frac{6\,(2 n - 3)}{n + 1}
\;=\; 12 \;-\; \frac{30}{n + 1},
\]
giving $C^{W{\rm-line}}_{\cW_3} = 12 = 2 \cdot 6$. The doubling
factor $2$ decomposes as the product of the Hessian ratio $3/2$ and
the kernel ratio $4/3$ at the $W$-current insertion. Stirling
asymptotic: $a_n \sim 12^{n}/(72\,\sqrt{\pi}\,n^{5/2})$. The lane is
Koszul-dual invariant under $c \mapsto 4 - c$ at conductor $K_3 = 4$.
\end{proposition}

\begin{remark}[Spin-stratified lattice, conjectural extension]
\label{rem:spin-stratified-cwsl}
The principal $\cW_N$ algebra has a sequence of primary
generators of conformal weights $s = 2, 3, \dots, N$. Conjectural
spin-stratified lattice:
\[
C_{\cW^{(s)}} \;=\; s\,(s + 1)
\qquad
\text{for the spin-$s$ primary in principal } \cW_N.
\]
At $s = 2$: $C = 6$ (Virasoro $T$-line, proved). At $s = 3$:
$C = 12$ ($\cW_3$ $W$-line, proved). At $s = 4$:
$C = 20$ (conjectural, $\cW_4$).
\end{remark}

\begin{proposition}[Cyclotomic singularities $\omega \in \mathbb Q(\zeta_6)$, iter~$37$]
\label{prop:cyclotomic-singularities-iter37}
\ClaimStatusProvedElsewhere
The Borel-plane singularities $\omega = 1/t_\pm(c)$ of the Virasoro
shadow lie in the cyclotomic field $\mathbb Q(\zeta_6)$. The discrete
$\mathbb Z/6$-symmetry is the dihedral symmetry of the Hessian
factorisation of $Q_c$.
\end{proposition}

\begin{proposition}[Lee--Yang phase at $c = -22/5$, iter~$46$]
\label{prop:lee-yang-phase}
\ClaimStatusProvedElsewhere
At $c = -22/5$ (the Lee--Yang minimal model central charge), the
Zamolodchikov factor $5 c + 22$ vanishes, and the Riccati
discriminant $\Delta_c = 40/(5 c + 22)$ diverges. The associated
shadow phase is the Lee--Yang singularity at the boundary of the
unitary regime; the $S_4$ pole signals a phase transition in the
shadow tower.
\end{proposition}

\begin{proposition}[Double-root phase at $c = -83/20$, iter~$44$--$48$]
\label{prop:double-root-phase}
\ClaimStatusProvedElsewhere
The Riccati polynomial $Q_c(t)$ acquires a double root at
$c = -83/20$, where the discriminant
$\disc_t Q_c = -8 \kappa^{2}\,\Delta_c$ vanishes by an arithmetic
coincidence of the rational map $c \mapsto (\kappa(c), \alpha,
S_4(c))$. The spectral curve $\Sigma_c$ acquires a node at this $c$;
the hyperelliptic involution degenerates to the identity on the
nodal fibre, and the shadow tower transitions into a degenerate
class L-like regime locally at $c = -83/20$.
\end{proposition}

\begin{proposition}[Large-$c$ expansion of $|\omega|(c)$, iter~$35$]
\label{prop:omega-large-c-expansion}
\ClaimStatusProvedElsewhere
The modulus of the leading Borel singularity admits the large-$c$
expansion
\[
\lvert\omega\rvert(c) \;=\; 2 \;-\; \frac{1}{4\,c}
\;+\; O\!\bigl(1/c^{2}\bigr)
\qquad (c \to \infty),
\]
with leading constant $2$ matching the asymptotic
$|t_\pm| \to 1/2$ as $c \to \infty$ (where the Riccati polynomial
$Q_c \to (2\kappa)^{2} + 12\kappa\,t + 9\,t^{2}$ in normalised
form), and subleading $-1/(4 c)$ encoding the leading
finite-central-charge correction.
\end{proposition}

\begin{proposition}[$\beta_N = 12 (H_N - 1)$ per-spin lane, Vol II inscription]
\label{prop:beta-N-per-spin-lane}
\ClaimStatusProvedElsewhere
The Vol II inscription
$\beta_N = \sum_{s = 2}^{N} 12/s = 12\,(H_N - 1)$ records the
per-spin-$s$ contribution to the principal $\cW_N$ tempering rate.
The contribution at spin $s$ is $12/s$. Boundary check at $N = 2$
gives $\beta_2 = 12 \cdot 1/2 = 6$, matching the Virasoro
$T$-line $C_\cA = 6$ of
Proposition~\ref{prop:universal-base-CA-six}. The closed form
$12 (H_N - 1) = 12 \sum_{s = 2}^{N} 1/s$ is rational (not integer)
for $N \ge 5$, with $\beta_5 = 77/5$ and $\beta_6 = 87/5$. The
closed form supersedes both prior candidates $(N+1)(N+2)/2$
(predicts $15$ at $N = 4$) and $N^{2} - N + 4$ (predicts $16$ at
$N = 4$); the true value is $\beta_4 = 13$ matching $12\,(1/2 +
1/3 + 1/4) = 12 \cdot 13/12 = 13$.
\end{proposition}

\begin{proposition}[$\cW(p)$ triplet $T$-line / Cartan-line shadow engines, iter~$77$]
\label{prop:wp-triplet-T-Cartan-line}
\ClaimStatusProvedElsewhere
For the logarithmic triplet algebra $\cW(p)$ at $c = 1 - 6 (p -
1)^{2}/p$, two parallel shadow lanes have been engineered:
\begin{enumerate}[label=\textup{(\roman*)}]
\item the $T$-line (stress-tensor lane), with shadow coefficients
$S_r(\cW(p)|_T)$ computed for $r = 2, \dots, 10$ at $p = 2$ and at
$p = 3$;
\item the Cartan-line (triplet lane), with shadow coefficients
$S_r(\cW(p)|_{\rm Cartan})$ computed for $r = 2, \dots, 10$ at
$p = 3$, including a full-scaffold quartic match at iter~$77$.
\end{enumerate}
The two lanes carry distinct asymptotic constants and exhibit the
expected Massey-product growth pattern; the structural conclusion
that $\cW(p)$ tempers via Zhu-bounded Massey is RETRACTED at
publication standard (see \S\ref{sec:open-frontier}).
\end{proposition}

\section{Per-family witnesses}
\label{sec:per-family-witnesses}

\paragraph{Heisenberg $\Heis_k$ (class G).}
Modular characteristic $\kappa(\Heis_k) = k$ on the current line
(landscape census, family G). Cubic $\alpha = 0$; quartic $S_4 = 0$;
discriminant $\Delta = 0$. The Riccati polynomial collapses to
$Q_c \equiv (2 k)^{2}$, a constant. Shadow tower terminates at $r =
2$: $H(t) = 2 k\,t^{2}$, the unison voice. Free-field bar; CG opening
of the chapter.

\paragraph{Affine Kac--Moody $V_k(\fg)$, $\fg$ simple, $k \neq -\hv$ (class L).}
Modular characteristic $\kappa(V_k(\fg)) = \dim(\fg)\,(k + \hv) /
(2\,\hv)$ (census C3). Cubic $\alpha = 1$ from the Killing form
(structure-constant $3$-cocycle). Quartic $S_4 = 0$; discriminant
$\Delta = 0$. The Riccati polynomial is a perfect square with linear
remainder: $Q_c = (2\kappa + 3 t)^{2}$. Shadow tower terminates at
$r = 3$: $H(t) = t^{2}(2\kappa + 3 t)$. Critical level $k = -\hv$
excluded by the Sugawara-degeneration caveat of Theorem C3
(landscape census, kappa vanishes at $k = -\hv$).

\paragraph{$\beta\gamma$ ghost system (class C).}
The class-C witness is the standard conformal-weight family
$\beta\gamma_\lambda$:
\[
S_2 = 6\lambda^2 - 6\lambda + 1,\qquad
S_3 = 0,\qquad
S_4 = -5/12,\qquad
S_r = 0 \quad (r \ge 5).
\]
Hence $r_{\max} = 4$ and $d_{\rm alg} = 2$. The internal
weight-changing line has zero shadow tower by rank-one abelian
rigidity, while the T-line is class~M. The class-C witness is the
full standard $\lambda$-family, not either one-dimensional slice.

\paragraph{Virasoro $\Vir_c$ (class M).}
Modular characteristic $\kappa(\Vir_c) = c/2$ (census C3). Cubic
$\alpha = 2$. Quartic $S_4 = 10/[c\,(5 c + 22)]$ (census). Discriminant
$\Delta = 8 \kappa\,S_4 = 40/(5 c + 22) \neq 0$ for generic $c$. The
Riccati polynomial is irreducible:
$Q_c = (c + 6\,t)^{2} + 80\,t^{2}/(5 c + 22)$. Shadow tower of depth
$\infty$; spectral curve $\Sigma_c$ smooth; Borel summable on the
divergent regime $c < c^{*}$, with $c^{*} \approx 6.125$ the unique
positive real root of $5 c^{3} + 22 c^{2} - 180 c - 872 = 0$.

\paragraph{Principal $\cW_3$ (class M).}
Modular characteristic $\kappa(\cW_3) = 5 c/6$ (census C4 at
$N = 3$). Two primary lanes: the $T$-line (Virasoro $C_\cA = 6$) and
the $W$-line ($C^{W{\rm-line}}_{\cW_3} = 12$,
Proposition~\ref{prop:w3-Wline-twelve}). Both lanes are class M;
the spectral curves $\Sigma_c^{(T)}$ and $\Sigma_c^{(W)}$ are
distinct, with distinct Borel singularities.

\paragraph{Bershadsky--Polyakov $\BP_k$ (class M).}
Convention-dependent central-charge formula. In the Arakawa
convention $c(\BP_k) = 2 - 24 (k + 1)^{2}/(k + 3)$, the Koszul
conductor $K_{\BP}(k) = c(\BP_k) + c(\BP_{-k - 6}) \equiv 196$ as a
constant rational function of $k$, with the fixed point $k = -3$
coinciding with the critical level $-\hv(\fsl_3)$.
$\kappa(\BP_k) + \kappa(\BP_k^{!}) = 98/3$ (self-dual at $k = -3$).
Class M with depth $\infty$. (Vol I standalone
\texttt{bp\_self\_duality.tex} carries the polynomial-identity
theorem in Arakawa convention; the Vol II Fateev--Lukyanov
convention $c^{\rm FL}(k) = -(2 k + 3)(3 k + 1)/(k + 3)$ has the
same algebra but produces meromorphic rather than polynomial
$K^{\rm FL}_{\BP}$. Both conventions parametrise the same algebra by
level reparametrisation.)

\paragraph{Lattice VOA $V_\Lambda$ (class G).}
Modular characteristic $\kappa(V_\Lambda) = \mathrm{rank}(\Lambda)$
(landscape census, lattice family). For Niemeier lattices,
$\kappa = 24$ blind to the root system. Cubic $\alpha = 0$; quartic
$S_4 = 0$; discriminant $\Delta = 0$. Class G uniformly.

\section{Open frontier}
\label{sec:open-frontier}

The chapter records the proved core. We catalogue the two open
fronts that survive the 2026-04-17 Beilinson rectification.

\paragraph{Logarithmic $\cW(p)$ tempering (RETRACTED to Conjectured).}
The Vol II commit \texttt{a5640de} retracts
\texttt{conj:tempered-stratum-contains-wp} from
\textsc{ProvedHere} to \textsc{Conjectured}: the proof chain
``$C_2$-cofiniteness $\Rightarrow$ bounded Massey $\Rightarrow$
tempered'' fails for $\cW(p)$. Gurarie 1993 (arXiv:hep-th/9303160)
and Flohr 1996 (arXiv:hep-th/9605151) construct logarithmic-CFT
amplitudes on $\cW(p)$ with unbounded Massey products despite
finite-dimensional Zhu algebra. The correct tempered scope is:
principal, non-logarithmic, non-minimal standard landscape; $\cW(p)$
is open pending the Adamovi\'c--Milas character-amplitude bound. The
$T$-line and Cartan-line shadow engines of
Proposition~\ref{prop:wp-triplet-T-Cartan-line} provide numerical
data through $r = 10$ at $p = 2, 3$, but the structural tempering
claim is downgraded.

\paragraph{Original-complex chain-level for class M.}
The chain-level MC5 statement on the original (non-weight-completed)
bar complex of class M chirally Koszul vertex algebras is
genuinely \emph{false} at the direct-sum level, since
$S_4(\Vir_c) = 10/[c\,(5 c + 22)] \neq 0$ at generic $c$ forces
non-trivial bar cohomology in weight $4$. The correct scope for the
chain-level statement is the canonical pro-object / $J$-adic /
filtered weight-completed ambient; the proof is now
\textup{(}Theorems~\ref{thm:mc5-class-m-chain-level-pro-ambient},
\ref{thm:mc5-class-m-topological-chain-level-j-adic};
Proposition~\ref{prop:ambient-equivalence}\textup{)}. The bounded
direct-sum statement is not a goal but a category
error; the resolution is to work in the weight-completed
category, where the Milnor / Mittag-Leffler passage to the inverse
limit absorbs the harmonic discrepancies at each weight stage.

\section{The single Platonic display}
\label{sec:platonic-display}

The chapter opens and closes with one identity. The four classes
G, L, C, M are the four factorisation types of $Q_c$. The shadow
connection is the Picard--Fuchs of $\Sigma_c = \{y^{2} = Q_c\}$. The
asymptotic resummation is by Pad\'e where the series converges and
by Borel where it diverges, with the crossover at the critical cubic
$5 c^{3} + 22 c^{2} - 180 c - 872 = 0$ for Virasoro. The
trans-series is the alien calculus of $\widehat S$ around its
branch points. The mock modularity is conjecturally the Borcherds
lift of $Q_c$ on $\Sigma_c$.

\begin{equation}
\label{eq:platonic-display}
\boxed{\;
H(t)^{2} \;=\; t^{4}\,Q_c(t),
\qquad
Q_c(t) \;=\; (2\kappa + 3\alpha\,t)^{2} \;+\; 2\,(8\,\kappa\,S_4)\,t^{2}.
\;}
\end{equation}

\begin{summary}[The chapter in one paragraph]
\label{sum:chapter-one-paragraph}
The shadow tower
$\{S_r(\cA)\}_{r \ge 2}$ on a primary line $L$ of a chirally Koszul
vertex algebra of finite type is governed by the algebraic identity
$H^{2} = t^{4}\,Q_c$, with $Q_c$ a quadratic polynomial in three
genus-zero data $(\kappa, \alpha, S_4)$ read off from the OPE at
the level of the $2$-, $3$-, $4$-point function. The four classes
G, L, C, M of the standard landscape are the four factorisation
types of $Q_c$, in codimensions $2, 1, 1, 0$ of the parameter space.
Class M acquires a hyperelliptic spectral curve
$\Sigma_c = \{y^{2} = Q_c\}$ whose periods are the flat sections of
the Picard--Fuchs shadow connection $\nabla^{\rm sh}$. Five
corollaries follow: the MC element $\Theta_\cA$ satisfies the
Riccati identity (C1); Picard--Fuchs propagates the dependence
$(\kappa, \alpha, S_4)$ across the genus tower (C2); the Gevrey-$1$
divergent series is Borel-summable on the divergent regime, with
Pad\'e the correct tool on the convergent regime (C3); the
shadow--Feynman dictionary $F_L = \Sh_{L+1}$ is a Getzler--Kapranov
operadic boundary identity (C4); the bar Euler characters are
conjecturally mock modular by Borcherds lift (C5). The Borel-geometric
layer (iter~$30$--$77$) provides the explicit constants:
$C_\cA = 6$ on the non-logarithmic class M $T$-line,
$C^{W{\rm-line}}_{\cW_3} = 12 = 2 \cdot 6$ on the $\cW_3$ $W$-line,
cyclotomic singularities $\omega \in \mathbb Q(\zeta_6)$, large-$c$
expansion $|\omega|(c) = 2 - 1/(4 c)$, and the per-spin tempering
rate $\beta_N = 12 (H_N - 1)$ on principal $\cW_N$.
\end{summary}

\section{Synthesis}
\label{sec:stqp-supplement}

\begin{remark}[Class-$\mathcal M$ Virasoro face as the shadow-tower quadrichotomy apex]
\label{rem:stqp-1}
The shadow-tower quadrichotomy -- Mathieu class $\mathcal M$ vs Conway
class $\mathcal C$ vs umbral class $\mathcal U$ vs moonshine class $\mathcal S$ --
meets at the K3 base point through the class-$\mathcal M$ Virasoro face,
with shadow tower $m_3 = ?, m_5 = -80/9, m_7 = -24320/81, m_8 = 33157760/19683$
and $G_5^{\mathrm{conn}} = 775/5184$ (from direct computation). The
class-$\mathcal S$ face $c_{2d} = -214$ sits below, with umbral-moonshine
shift $m_\sigma$ (Cheng--Duncan--Harvey 2014) quantising the shift
between $\mathcal M$ and the other three classes \cite{ChengDuncanHarvey2014,EguchiOoguriTachikawa2010}.
Cross-ref Vol.~III Chapter~\ref{ch:k3-chiral-bialgebra-platonic}.
\end{remark}

\begin{remark}[$\mathbf H_{\Delta_5}$ as unified apex of the four shadow faces]
\label{rem:stqp-2}
The fourth taxon $\mathbf H_{\Delta_5}$ is the unified apex of the
four shadow faces: (i) Mathieu $\mathcal M_{24}$-twined trace gives the
Igusa $\Delta_5$ as frame shape (conjectural at d=2 -- see Vol.~III
mock modular theorem); (ii) Conway $\mathcal C$ shadow via $II_{1,1}\oplus II_{8}\oplus II_{8}\oplus II_{8}$
Leech embedding; (iii) umbral $\mathcal U$ shadow via $\widetilde M_{24}$-lift
obstruction; (iv) $\mathcal S$-class shadow via class-$\mathcal S$ $A_1$
on $\Sigma_{0,24}$ with $c_{2d} = -214$. The shadow-tower $m_k$
coefficients encode the $A_\infty$ coproduct corrections
\cite{ConwaySloane1988,ChengDuncanHarvey2014,Gaiotto2009}.
Cross-ref Chapter~\ref{ch:k3-chiral-bialgebra-platonic}.
\end{remark}

\section{K3 shadow spectrum at $c_{2d} = -214$}
\label{sec:stqp-supplement-ii}

\begin{remark}[Retraction of the $c_{2d} = -312$ identification]
\label{rem:stqp-retract-312}
\index{retraction!$c_{2d}=-312$ to $c_{2d}=-214$}
An earlier attempt read the class-$\mathcal S$ $A_1$ chiral
central charge on $\Sigma_{0,24}$ with all maximal punctures as
$c_{2d} = -312 = -12\cdot 26 = -\chi(\mathrm{K3})\cdot S_2(\Vir_{26})$
is withdrawn. The 4d conformal anomaly $c_{4d} = 26$ does not arise
from the Shapere--Tachikawa normalisation
$c_{4d} = (2n_v + n_h)/12$ applied to the class-$\mathcal S$ count
$(n_v, n_h) = (63, 88)$ for $(A_1, \Sigma_{0,24})$: direct
substitution yields $c_{4d} = (2\cdot 63 + 88)/12 = 214/12 = 107/6$,
not $26$. The Beem--Rastelli factor $c_{2d} = -12 c_{4d}$
(Beem--Lemos--Liendo--Peelaers--Rastelli--van~Rees 2013, eq.~(3.14))
then gives $c_{2d} = -214$, not $-312$. The correct first-principles
derivation is given in Vol~I Chapter~\ref{ch:thqg-open-closed-realization},
Proposition at line~1890 (Shapere--Tachikawa 2008, \texttt{arXiv:0804.1957};
Chacaltana--Distler 2010 eq.~(5.14)). The surrounding shadow-tower
propositions and mock-modular weight statements below are corrected
to the true value $c_{2d} = -214$; every numerical invariant is
recomputed. The six-fold pentagon-umbral $\mathbb{Z}/6$ anomaly
and the class-$\mathrm M$ placement of $\mathbf H_{\Delta_5}$ survive
unchanged; only the central-charge arithmetic changes.
\end{remark}

The Shapere--Tachikawa class-$\mathcal S$ normalisation
derives $c_{2d} = -214$ from first principles and decomposes
$-214 = -2\cdot 107$ into a shadow-tower spectrum on class~$\mathrm M$,
anchoring the apex of the Virasoro archetype at the K3 base point.

\begin{proposition}[$c_{2d} = -214$ shadow-tower decomposition;
\ClaimStatusProvedHere]
\label{prop:stqp-312-factor}
\index{central charge!$-214$!shadow-tower decomposition}
\index{K3!shadow tower class $\mathrm M$}
The class-$\mathcal S$ $A_1$ conformal anomaly on $\Sigma_{0,24}$
with $24$ maximal regular $\mathfrak{su}(2)$ punctures evaluates to
\[
c_{4d}(A_1, \Sigma_{0,24}, \mathrm{all\ max})
\;=\;
\frac{5n - 13}{6}\bigg|_{n=24}
\;=\;
\frac{5\cdot 24 - 13}{6}
\;=\;
\frac{107}{6},
\]
via the Shapere--Tachikawa normalisation
$c_{4d} = (2n_v + n_h)/12$ with class-$\mathcal S$ count
$(n_v, n_h) = (63, 88)$ for $(A_1, \Sigma_{0,24})$
(Shapere--Tachikawa 2008, arXiv:0804.1957; Chacaltana--Distler 2010
eq.~(5.14); Vol~I Chapter~\ref{ch:thqg-open-closed-realization}
line~1890). The integer $107$ is prime, so the Beem--Rastelli
image
\[
c_{2d}(A_1, \Sigma_{0,24}, \mathrm{all\ max})
\;=\;
-12\, c_{4d}
\;=\;
-2\cdot 107
\;=\;
-214
\]
admits no non-trivial integer factorisation apart from $-2\cdot 107$
(Beem--Lemos--Liendo--Peelaers--Rastelli--van~Rees 2013, eq.~(3.14)).
The linear structure $5n - 13$ reads as $5n$ vector-contribution
minus the universal Argyres--Seiberg gravitino subtraction $13$
(Shapere--Tachikawa 2008, \S 3).
Substituting $c = -214$ into the universal Virasoro shadow
tower of Proposition~\ref{prop:vir-shadow-r5}:
\[
S_2 \;=\; \frac{c}{2} \;=\; -107,
\qquad
S_3 \;=\; 2,
\]
\[
5c + 22 \;=\; -1048,
\qquad
c(5c+22) \;=\; (-214)\cdot(-1048) \;=\; 224\,272,
\]
\[
S_4 \;=\; \frac{10}{c(5c+22)} \;=\; \frac{10}{224\,272} \;=\; \frac{5}{112\,136},
\]
\[
S_5 \;=\; \frac{-48}{c^2(5c+22)} \;=\; \frac{-48}{(-214)^2\cdot(-1048)} \;=\; \frac{3}{3\,000\,319}
\]
(with $(-214)^2\cdot(-1048) = -48\,005\,104$;
$S_5 = -48/(-48\,005\,104) = 3/3\,000\,319$ after common-factor
cancellation). All three invariants $(S_2, S_4, S_5)$ are nonzero
and finite at $c = -214$; the shadow tower has depth
$r_{\max} = \infty$, and the K3 chiral bialgebra
$\mathbf H_{\Delta_5}$ inherits class~$\mathrm M$ (Virasoro
archetype) on the stress-tensor line.
\end{proposition}

\begin{proof}
The 4d count $(n_v, n_h) = (63, 88)$ on $(A_1, \Sigma_{0,24})$ with
all maximal punctures is obtained via the class-$\mathcal S$
polynomial $n_v(T_N) = (N-1)(N-2)(N+2)/2$, $n_h(T_N) = 2N(N^2-1)/3$
at $N = 2$ (giving $(0, 4)$ per $T_2$-trinion) summed over the
$22$ pants decomposition with $21$ gauge tubes carrying
$\mathrm{SU}(2)$ vectors $n_v = N^2-1 = 3$ each: $n_v = 21\cdot 3
+ 22\cdot 0 = 63$, $n_h = 22\cdot 4 = 88$
(Chacaltana--Distler 2010 eq.~(5.14); Vol~I
Chapter~\ref{ch:thqg-open-closed-realization}, Proposition
around line~1890). Shapere--Tachikawa 2008 arXiv:0804.1957
normalisation
$c_{4d} = (2n_v + n_h)/12 = (126 + 88)/12 = 214/12 = 107/6$.
Cross-check: the formula reduces to $c_{4d}(n) = (5n - 13)/6$
for the all-max sphere, yielding at $n = 3$ the free trinion
value $(a, c) = (1/6, 1/3)$ (so $c_{4d} = (15-13)/6 = 1/3$),
at $n = 4$ the $\mathrm{SU}(2)$ $N_f = 4$ SQCD value
$c_{4d} = 7/6$ (Tachikawa 2013 lectures eq.~(13.6)), at
$n = 24$ the K3 value $107/6$. Beem--Rastelli
$c_{2d} = -12 c_{4d} = -214$ (BLLPRvR 2013 eq.~(3.14)).

The shadow-tower values at $c = -214$ follow by direct substitution
into Proposition~\ref{prop:vir-shadow-r5}:
$5c + 22 = -1048$, $c(5c+22) = 224\,272 = 2^4\cdot 107^2\cdot\ldots$
(explicit factorisation: $224\,272 = 16\cdot 14\,017
= 2^4\cdot 14\,017$; the number $14\,017 = 107\cdot 131$ is a
product of two primes, neither of which enters Virasoro
minimal-model locus $5c+22 = 0$). The Zamolodchikov norm
$\langle\Lambda|\Lambda\rangle = c(5c+22)/10 = 224\,272/10
= 112\,136/5 = 22\,427.2 > 0$ is positive, since both factors
$c = -214$ and $5c+22 = -1048$ are negative; the
$\Lambda = {:}TT{:} - (3/10)\partial^2 T$ quasi-primary has
positive Zamolodchikov norm, confirming that $c = -214$ lies
outside the minimal-model degeneracy locus. The Riccati
polynomial $Q_c = (c + 6t)^2 + 80t^2/(5c+22)$ of
\eqref{eq:platonic-display} has nonvanishing discriminant
$\Delta_{\mathrm{Ricc}} = 40/(5c+22) = 40/(-1048) = -5/131
\neq 0$, placing $c = -214$ in the Borel-summable stratum of
the critical cubic $5c^3 + 22c^2 - 180c - 872$. Class~$\mathrm M$
(Virasoro archetype, $r_{\max} = \infty$) follows.
\end{proof}

\begin{remark}[Unitarity scope of $c = -214$]
\label{rem:stqp-unitarity-scope}
The central charge $c_{2d} = -214$ is \emph{not unitary} in the
positive-definite Hilbert-space sense: the Kac determinant
$\det G_h = \prod(h - h_{p,q}(c))^{P(h-pq)}$ has negative values at
positive $h$ for any $c < 0$, including $c = -214$. This is the
expected reading: $-214$ is the 2d chiral anomaly matching the
4d Weyl anomaly $c_{4d} = 107/6$ via Beem--Rastelli, descended
from the 6d $(2,0)$ $A_1$ anomaly polynomial integrated over
$\mathrm{K3}\times\Sigma_{0,24}$ (Alday--Benini--Tachikawa 2009,
arXiv:0912.4664; Shapere--Tachikawa 2008 \S 3), not a unitary
Hilbert-space central charge. The class-$\mathrm M$ shadow tower
of Proposition~\ref{prop:stqp-312-factor} is therefore
read in the \emph{effective conformal anomaly} sense
$c_{\mathrm{eff}} = c - 24 h_{\min}$, with $h_{\min}$ the
lowest Virasoro eigenvalue of the BKM BPS Fock module on
$\mathbf H_{\Delta_5}$. See Vol~II
Remark~\ref{rem:anomtop-312-spectrum} for the
anomaly-polynomial derivation at $c_{2d} = -214$ and the BPS
counting.
\end{remark}

\begin{remark}[Mock-modular shadow weight from $c = -214$;
Siegel weight $5$]
\label{rem:stqp-mock-weight}
\index{mock modular!Siegel weight 5}
\index{BKM denominator!weight!$\Delta_5$}
The naive weight reading $-c_{2d}/24 = 214/24 = 107/12$ is
\emph{not} integer, so the Borcherds lift of the K3 elliptic genus
cannot be packaged as a weight-$-c/24$ modular form on
$\mathrm{SL}_2(\mathbb Z)$. The correct modular-weight assignment
passes through the Gritsenko--Nikulin 1996--98 Siegel lift: the
$\mathrm T^2$ chiral partition function of the K3 BKM datum is a
Siegel modular form of weight~$5$ on $\mathrm{Sp}_4(\mathbb Z)$
(or its double cover carrying the half-integral Gritsenko-lift
character), identified with $\Delta_5 = \Phi_{10}^{1/2}$ through
the Borcherds denominator identity
$\prod_{r\cdot\alpha > 0}(1 - q^{r\alpha})^{c(r\alpha)}
= \Phi_{10}(Z)$ on $\mathrm{Sp}_4(\mathbb Z)\backslash\mathbb H_2$
(Borcherds 1992 Theorem~10.1; Gritsenko--Nikulin 1998
\S 4; Cheng--Duncan--Harvey 2014). The Siegel weight
$w_{\mathrm{Sieg}} = 5$ is determined by the Borcherds
denominator degree $10$ and the square-root refinement
$\Delta_5 = \sqrt{\Phi_{10}}$, independent of the elliptic
coordinate $c_{2d}$. The bridge from Borcherds-Siegel weight to
elliptic-fibre weight is the Gritsenko--Nikulin restriction to
the Humbert divisor $H_1$, under which $\Delta_5$ descends to a
weight-$5$ Jacobi form; this is distinct from the naive
Cardy-regime expression $-c_{2d}/24$. Mock completion follows
Zwegers 2002, Cheng--Duncan--Harvey 2014, and
Bringmann--Creutzig--Rolen 2014 on K3 BKM partition functions.
The non-integrality of $107/12$ is thus a feature, not a bug: it
signals that the correct modular covariance is Siegel rather than
elliptic, consistent with the K3 base point being governed by
$\mathrm{Sp}_4(\mathbb Z)$ via Kuga--Satake rather than by
$\mathrm{SL}_2(\mathbb Z)$ alone.
\end{remark}

\begin{remark}[Cardy counting at $c = -214$]
\label{rem:stqp-cardy}
\index{Cardy formula!effective central charge}
\index{BPS counting!K3 BKM}
The Cardy asymptotic density $S(E) = 2\pi\sqrt{c E/6}$ is ill-defined
at $c = -214 < 0$. The physically correct count uses the effective
central charge $c_{\mathrm{eff}} = c - 24 h_{\min}$, with
$h_{\min}$ the lowest Virasoro eigenvalue in the BKM BPS module,
equivalently the polar part of the Borcherds denominator identity
at the cusp. For $\mathbf H_{\Delta_5}$ the polar part of
$1/\Phi_{10}(Z)$ has weight $-10$ at the maximal-rank cusp of
$\mathcal A_2$, giving $h_{\min} = -2$ (half the Borcherds
denominator weight, after the Gritsenko--Nikulin double-cover
normalisation), so
\[
c_{\mathrm{eff}} \;=\; -214 - 24(-2) \;=\; -214 + 48 \;=\; -166.
\]
This remains negative, consistent with the gravitational-anomaly
sign reading of Remark~\ref{rem:stqp-unitarity-scope}:
BKM BPS states are counted against an anti-holomorphic background,
and the Cardy density is meaningful on the modulus-squared
$|Z|^2$ rather than the bare holomorphic part.
See Witten 2007, Manschot--Moore 2007, Dijkgraaf--Verlinde--Verlinde
1997 for the $|c_{\mathrm{eff}}|$-based counting; the mapping to
K3 BPS invariants uses the Harvey--Moore 1996 singular-theta
correspondence. Note that the effective value
$c_{\mathrm{eff}} = -166 = -2\cdot 83$ (with $83$ prime) at
$c_{2d} = -214$ differs numerically from the retracted
value $c_{\mathrm{eff}} = -264$ at $c_{2d} = -312$; the
mathematical content of the Borcherds-polar-part identification
is unchanged.
\end{remark}

\begin{remark}[Class-$\mathrm M$ placement of K3 via $c = -214$]
\label{rem:stqp-class-m}
\index{K3!class $\mathrm M$ placement}
\index{shadow tower!K3 depth-infty}
Combining Proposition~\ref{prop:stqp-312-factor} with the
class-$\mathrm M$ Riccati criterion of
Section~\ref{sec:class-M-Riemann-surface}: the K3 BKM chiral
stress-tensor line has
$(\kappa, \alpha, S_4) = (-107, 2, 5/112\,136)$ with discriminant
$\Delta_{\mathrm{Ricc}} = 8\kappa S_4 = -10/131 \neq 0$, so $Q_c$
is irreducible and $r_{\max} = \infty$. The hyperelliptic spectral
surface $\Sigma_{c = -214} = \{y^2 = Q_{-214}(t)\}$ has smooth
branch structure away from $t = 0$ and
$t = -2\kappa/(3\alpha) = 107/3$, governed by Picard--Fuchs as in
Theorem~\ref{thm:picard-fuchs-shadow}. Hence the K3 chiral
bialgebra $\mathbf H_{\Delta_5}$ is firmly class~$\mathrm M$
(Virasoro archetype with $r_{\max} = \infty$) on its stress-tensor
line; it is not class~$\mathrm G$ (lattice VOAs have
$\alpha = S_4 = 0$, Section~\ref{sec:class-M-Riemann-surface}
contrasting paragraph on lattice VOAs). The fourth taxon
$\mathbf H_{\Delta_5}$ of Remark~\ref{rem:stqp-2} is a
class-$\mathrm M$ object refined by $M_{24}$-umbral shadow and
Kuga--Satake geometry rather than a fifth class. The class-$\mathrm M$
conclusion is \emph{robust under the $-312 \to -214$ retraction}:
both central-charge values lie in the generic class-$\mathrm M$
stratum (nonvanishing $5c + 22$, nonvanishing $S_4$,
irreducible $Q_c$); the depth-$\infty$ tower is a topological
invariant of the stratum, not of the specific central charge.
\end{remark}

\begin{remark}[Conformal-block basis at $c_{2d} = -214$ and the
crossing identity]
\label{rem:stqp-block-basis-crossing}
\index{conformal block!$\Delta_5$!$c=-214$!block count $130$}
\index{crossing symmetry!$\Delta_5$!Humbert $H_1$}
The class-$\mathrm M$ depth-$\infty$ tower at $c_{2d} = -214$
admits a finite genus-one conformal-block basis on the Humbert
divisor $H_1\subset\mathfrak H_2$: truncating the BKM positive
real roots of $\mathfrak g_{\Delta_5}$ at height $\ell = 8$ yields
a $130$-element block family
$\mathcal F_{\Delta_5, 8} = \{F_\lambda(\tau)\}$
indexed by $\Lambda^{+}_{\Delta_5, 8}\cup\{0\}$, with Fourier
coefficients extracted from $1/\Delta_5$ via the
Gritsenko--Nikulin 1998 \S 3.5 restriction to $\{z = 0\}$. The
block count $130 = 1 + 129$ matches the Vol~II PBW-finite bound
$8^{129}$ at $N = 129$ real positive roots (Vol~II
Theorem~\ref{V2-thm:modbootstrap-cross-block-count},
Vol~III Theorem~\ref{V3-thm:modtr-Plancherel-ML-limit}).
The modular-$S$ action is the Fricke involution $w_8$ acting by
$\lambda\mapsto -w_0(\lambda)$ on BKM weights (Vol~II
Theorem~\ref{V2-thm:modbootstrap-cross-Smatrix}), and the
torus partition sum
$Z_{\Delta_5}(\tau,\bar\tau) = \sum_\lambda |F_\lambda(\tau)|^2$
is $\mathrm{SL}_2(\mathbb Z)$-invariant
(Vol~II Proposition~\ref{V2-prop:modbootstrap-cross-Tmatrix}),
lifting to an $\mathrm{Sp}_4(\mathbb Z)$-invariant of weight
$10$ on the full $\mathcal A_2$
(Vol~II Theorem~\ref{V2-thm:modbootstrap-cross-214},
Remark~\ref{V2-rem:modbootstrap-Sp4}).
The Fricke-fixed subset
$\mathcal F^{w_8}_{\Delta_5, 8} =
\{F_0, F_\rho, F_{\rho + \alpha_{\mathrm{hyp}}}, F_{2\rho}\}$
has cardinality~$4$, matching the
$M_{24}$-invariant Humbert-intersection count
$|H_1\cap H_4|/M_{24} = 4$
(Vol~II
Theorem~\ref{V2-thm:modbootstrap-cross-four-fricke-fixed}).
On the non-semisimple locus
$(H_1\cup H_4)\setminus(H_1\cap H_4)$ (where the
$\mathbf H_{\Delta_5}$-MTC fails to be semisimple; the chiral
Koszul property is unaffected),
the non-unitary Virasoro sector at $c = -214$ acquires a
mock-modular shadow; the Zwegers 2002 \S 1.2 completion
$\widehat F_\lambda = F_\lambda + F_\lambda^*$ with
Cheng--Duncan--Harvey 2014 \S 4 umbral shadow restores
$\mathrm{SL}_2(\mathbb Z)$-covariance on the full
$\mathcal A_2$, and the crossing identity
$\sum_\lambda F_\lambda(\tau)\overline{F_\lambda(-1/\bar\tau)}
= \sum_\mu F_\mu(\tau+1)\overline{F_\mu(\overline{\tau+1})}$
lifts to the completed basis.
The primitive $12$th-root phase $e^{2\pi i\cdot 107/12}$ in the
$T$-action encodes the non-integrality of $-c/24 = 107/12$, which
signals Siegel-rather-than-elliptic modular covariance
(Remark~\ref{rem:stqp-mock-weight}); it cancels in $|F_\lambda|^2$
and leaves no obstruction to crossing symmetry on the Koszul
locus.
Primary: Gritsenko--Nikulin 1998 \S 3.5, \S 5.1, \S 6
(Humbert restriction, Weyl vector, Fricke fixed-point locus);
Belavin--Polyakov--Zamolodchikov 1984 \S 2 (crossing axiom);
Moore--Seiberg 1989 \S 2 ($S^2 = (ST)^3$);
Cheng--Duncan--Harvey 2014 \S 4, \S 6
(umbral shadow, $M_{24}$-twining table);
Zwegers 2002 \S 1.2, \S 3.1 (mock completion);
Conway--Sloane 1988 Ch.~10 (Niemeier-lattice decomposition of
$II_{2,1}\oplus II_8^{\oplus 3}$).
\end{remark}

\section{unitarity on the Koszul locus $\overline{\mathcal A_2}\setminus(H_1\cup H_4)$}
\label{sec:stqp-unitarity}

The BKM analysis sharpens the earlier scope
(Remark~\ref{rem:stqp-unitarity-scope}) by isolating \emph{which}
sub-module of the K3 BKM chiral bialgebra $\mathbf H_{\Delta_5}$ is
unitary, on the Koszul locus
$\overline{\mathcal A_2}\setminus(H_1\cup H_4)$. The ambient module is
indefinite-Hermitian; the positive-energy unitary submodule is the
real-root sector, identified through the first-principles signature
computation of the Feingold--Frenkel rank-$3$ hyperbolic Cartan.

\begin{proposition}[Hyperbolic Cartan signature $(2,1)$ for $\mathbf H_{\Delta_5}$;
\ClaimStatusProvedHere]
\label{prop:stqp-signature}
\index{hyperbolic Cartan!signature $(2,1)$}
\index{Feingold--Frenkel!rank-$3$ hyperbolic}
\index{BKM algebra!$\mathbf H_{\Delta_5}$!real-root Gram}
The rank-$3$ real simple roots $\alpha_1, \alpha_2, \alpha_3$ of
the BKM algebra $\mathfrak g_{\Delta_5}$ attached to the K3 BKM
chiral bialgebra $\mathbf H_{\Delta_5}$ span a hyperbolic
Cartan of signature $(2,1)$. The Gram matrix is
\[
G_{\mathrm{BKM}} \;=\;
\begin{pmatrix}
2 & -2 & -2 \\
-2 & 2 & -2 \\
-2 & -2 & 2
\end{pmatrix},
\qquad
\det G_{\mathrm{BKM}} \;=\; -32,
\]
with spectrum $\{4, 4, -2\}$: two positive eigenvalues $+4$
(positive-definite $2$-plane $P \subset \mathfrak h_{\mathbb R}$)
and one negative eigenvalue $-2$ (the hyperbolic direction
$L \subset \mathfrak h_{\mathbb R}$).
\end{proposition}

\begin{proof}
The Gram matrix is the Feingold--Frenkel rank-$3$ hyperbolic
Cartan matrix ($F$, Feingold--Frenkel 1983
arXiv:math/9210213 eq.~(1.1); Math.\ Ann.\ 263 (1983), 87--144,
Theorem~1.1), whose three real simple roots are identified with
the $(2,1)$-Lorentzian lattice $II_{1,1}\oplus\langle 2\rangle$
after Kac--Moody truncation. The characteristic polynomial is
$p(\lambda) = \det(G_{\mathrm{BKM}} - \lambda I)
= -\lambda^3 + 6\lambda^2 - 32$ (Vandermonde expansion on the
circulant-up-to-sign structure); its real roots are
$\lambda_1 = \lambda_2 = 4$ (double) and $\lambda_3 = -2$
(simple), so the signature is $(2, 1, 0)$: two positive, one
negative, zero null. The $\lambda = 4$ eigenspace is spanned
by the weight-$2$ combinations
$(\alpha_1 - \alpha_2)$ and $(\alpha_1 + \alpha_2 - 2\alpha_3)/\sqrt 3$
(orthonormal after normalisation), both of norm-squared
$4\cdot 2 = 8$ in unnormalised form; the $\lambda = -2$
eigenspace is spanned by the totally symmetric combination
$\alpha_+ = (\alpha_1 + \alpha_2 + \alpha_3)/\sqrt 3$ with
$(\alpha_+, \alpha_+) = -2\cdot 2 = -4$, identified with
the Lorentzian imaginary root direction. The double
eigenvalue $+4$ reflects the $S_3$ root-permutation
symmetry of the Feingold--Frenkel Dynkin diagram
(three vertices, all pairwise joined by bold edges
$a_{ij} = -2$), which acts irreducibly on the $\lambda = 4$
plane. The principal minor $m_2 = \det\bigl(
\begin{smallmatrix} 2 & -2 \\ -2 & 2
\end{smallmatrix}\bigr) = 0$ (not $4$ as one might expect from
simply-laced $A_2$): the $\{\alpha_1, \alpha_2\}$ sublattice has
isotropic combination $\alpha_1 + \alpha_2$, so the principal-minor
ladder is $(2, 0, -32)$ rather than $(2, 4, \det)$, and signature
must be read from eigenvalues rather than Sylvester. This
confirms the $(2, 1)$ signature is genuinely hyperbolic, not
semi-positive-definite as a hasty principal-minor read would
suggest. Independent verification via Gritsenko--Nikulin 1998
Theorem~1.1: the real-root lattice of
$\mathfrak g_{\Delta_5}$ is $II_{2,1}$ by construction
(paramodular weight~$5$ via singular-theta lift on the
orthogonal Grassmannian of signature $(2,1)$); so signature
$(2,1)$ is forced by the Borcherds lift data independent of
the Feingold--Frenkel Cartan computation, and the numerical
spectrum $\{4, 4, -2\}$ cross-verifies the theorem.
Epistemic hierarchy: direct eigenvalue computation ($+4, +4, -2$,
three times); characteristic-polynomial verification
$p(\lambda) = -\lambda^3 + 6\lambda^2 - 32$; Gritsenko--Nikulin
1998 Theorem~1.1 signature; Feingold--Frenkel 1983 Math.\ Ann.\
263 classification.
\end{proof}

\begin{remark}[Scope: signature $(2,1)$ vs naive $A_3$ or $D_3$]
\label{rem:stqp-signature-scope}
The hyperbolic signature $(2, 1)$ is what distinguishes the BKM
algebra $\mathfrak g_{\Delta_5}$ from a finite-dimensional
semisimple Kac--Moody algebra: the $(\alpha_1, \alpha_2) = -2$
entry (bold edge in the Dynkin diagram) rather than $-1$ (simple
edge) pushes the Cartan matrix out of the positive-definite
$A_3$-or-$D_3$ regime. A finite-type rank-$3$ simply-laced
Cartan would have eigenvalues $\{1, 3, 4\}$ ($A_3$) or
$\{2, 2, 4\}$ ($D_3$), all positive. The Feingold--Frenkel
bold-edge structure produces one negative eigenvalue, marking
$\mathfrak g_{\Delta_5}$ as genuinely infinite-dimensional
(Lorentzian Kac--Moody) with an infinite-dimensional root
lattice completed by imaginary roots (positive multiplicity
$c(n^2/2)$ from the $K3$ elliptic genus
$2\phi_{0,1}(\tau, z) = \sum c(n, \ell) q^n y^\ell$;
Eguchi--Ooguri--Tachikawa 2010 arXiv:1004.0956; Cheng 2010
arXiv:1005.5415). The $+4$ double eigenvalue is the
$E_8\oplus E_8$-like degeneracy preserved by the $S_3$ Weyl
symmetry; the $-2$ eigenvalue is the direction along which
imaginary roots accumulate at the maximal-rank cusp of
$\mathcal A_2$.
\end{remark}

\begin{proposition}[Positive-energy unitary spectrum; \ClaimStatusProvedHere]
\label{prop:stqp-unitary-spectrum}
\index{unitary representation!$\mathfrak g_{\Delta_5}$}
\index{positive-energy module!real-root sector}
\index{Frenkel--Kac!vertex algebra on lattice}
Let $P \subset \mathfrak h_{\mathbb R}$ denote the positive-definite
$2$-plane of Proposition~\ref{prop:stqp-signature}
(the $\lambda = 4$ eigenspace), and let $L$ denote the hyperbolic
direction (the $\lambda = -2$ eigenspace). A highest-weight
$\mathfrak g_{\Delta_5}$-module $V^\lambda$ with dominant integral
weight $\lambda \in \mathfrak h^*_{\mathbb Z}$ is positive-energy
unitary if and only if
\[
\lambda \in P^* \quad\text{and}\quad (\lambda, \lambda) > 0,
\]
where $P^*$ is the integral dual lattice of $P$ in $\mathfrak h^*$
and $(\cdot, \cdot)$ is the Gram form of
Proposition~\ref{prop:stqp-signature}.
Weights $\lambda$ with $(\lambda, \lambda) < 0$ (imaginary-root
sector, supported on $L$) yield non-unitary modules with
negative-norm states.
\end{proposition}

\begin{proof}
The representation theory of a BKM algebra over a Lorentzian
Cartan is the Frenkel--Kac construction on the positive-definite
projection: $V^\lambda = \mathrm{Fock}(P)\otimes\mathrm{Fock}(L)$
where the $P$-factor is a standard lattice vertex algebra on a
positive-definite rank-$2$ lattice (Frenkel--Kac 1980
Invent.\ Math.\ 62; Frenkel--Lepowsky--Meurman 1988
\S 8.8), and the $L$-factor is a rank-$1$ indefinite Fock
module. The $P$-factor has Hermitian inner product
$\langle\cdot,\cdot\rangle_P$ descending from the positive-definite
Gram on $P$ via the usual Heisenberg formula
$\langle e^\lambda, e^\mu\rangle = \delta_{\lambda,-\mu}$
with sign determined by $(\lambda, \lambda)_P > 0$; all
$L_0$-eigenvalues are non-negative, and the only null vector at
generic $(\lambda, \lambda)_P > 0$ is the highest-weight
vector itself, giving positive-energy unitarity on this factor.
The $L$-factor is a rank-$1$ indefinite Fock module on the
$(1,0)$ or $(0,1)$-Lorentzian signature, whose Zamolodchikov
norm $\langle\Lambda_L|\Lambda_L\rangle$ changes sign with
$(\lambda, \lambda)_L$; for $(\lambda, \lambda)_L > 0$ (vector,
after the total sign $(\lambda, \lambda) > 0$) the $L$-factor
admits a positive-definite contravariant form, and unitarity on
the tensor product is preserved. For $(\lambda, \lambda) < 0$
(imaginary roots, timelike projection dominant) the $L$-factor
carries a negative-norm ground state, and $V^\lambda$ is
non-unitary. This is the standard Lorentzian-lattice reading:
see Borcherds 1992 Invent.\ Math.\ 109 \S 4, Gebert--Nicolai
1995 arXiv:hep-th/9504153 \S 3, and Kac 1990 \emph{Infinite
Dimensional Lie Algebras} Chapter~14 on unitary highest-weight
modules over Lorentzian Kac--Moody. The positive-Weyl-chamber
restriction $\lambda \in P^*_+$ ensures the integrable condition
(Verma module quotient is a simple integrable module) via the
Garland--Lepowsky 1976 / Kac 1974 integrability criterion on
the real-root Weyl group $W(P)$; the strict positivity
$(\lambda, \lambda) > 0$ then forces $\lambda$ to be spacelike
for the Lorentzian form.
\end{proof}

\begin{corollary}[Real-root / imaginary-root dichotomy; \ClaimStatusProvedHere]
\label{cor:stqp-real-imag-dichotomy}
The unitarity spectrum of $\mathbf H_{\Delta_5}$ decomposes as
\[
\mathrm{Spec}^{\mathrm{unit}}(\mathbf H_{\Delta_5})
\;=\;
\{V^\lambda : \lambda \in \Phi^{\mathrm{re}}_+,\;
(\lambda, \lambda) > 0\},
\qquad
\mathrm{Spec}^{\mathrm{nonunit}}(\mathbf H_{\Delta_5})
\;=\;
\{V^\lambda : \lambda \in \Phi^{\mathrm{im}}_+,\;
(\lambda, \lambda) \le 0\},
\]
with $\Phi^{\mathrm{re}}_+$ the positive real-root set (finite
multiplicity~$1$) and $\Phi^{\mathrm{im}}_+$ the positive
imaginary-root set (multiplicity $c(n^2/2)$ from the K3 elliptic
genus $2\phi_{0,1}$). The indefinite Hermitian signature on the
BKM inner product is $(\#\Phi^{\mathrm{re}}, \#\Phi^{\mathrm{im}})
= (\infty, \infty)$: both sectors are infinite, and no finite
positive-definite projection captures the full BKM representation
content.
\end{corollary}

\begin{remark}[Effective central charge unitarity reading $c_{\mathrm{eff}} = -166$]
\label{rem:stqp-ceff-unitarity}
\index{effective central charge!$c_{\mathrm{eff}} = -166 = -2 \cdot 83$}
\index{non-unitary anomaly!effective $c$}
At the level of the bulk $c_{2d} = -214$, the Cardy-regime
effective central charge (Remark~\ref{rem:stqp-cardy}) is
\[
c_{\mathrm{eff}}
\;=\;
c_{2d} - 24\, h_{\min}
\;=\;
-214 + 48
\;=\;
-166
\;=\;
-2\cdot 83,
\]
with $83$ prime (so $-166$ admits no further nontrivial integer
factorisation). The unitary submodule of
Proposition~\ref{prop:stqp-unitary-spectrum} realises this
$c_{\mathrm{eff}} = -166$ as the \emph{effective} central charge
governing the finite-volume partition function of the real-root
sector: the character
\[
\chi^{\mathrm{unit}}_{\mathbf H_{\Delta_5}}(\tau)
\;=\;
\mathrm{tr}_{V^{\mathrm{unit}}} q^{L_0 - c_{\mathrm{eff}}/24}
\;=\;
\mathrm{tr}_{V^{\mathrm{unit}}} q^{L_0 + 83/12}
\]
converges on the upper half plane for $\mathrm{Im}(\tau) > 0$;
the negative sign of $c_{\mathrm{eff}}$ signals that
$\mathbf H_{\Delta_5}$ is not a unitary CFT at the whole-module
level, only on the real-root submodule. The $c_{\mathrm{eff}}/24
= -83/12$ shift appears in the Borcherds denominator identity
expansion as the weight of the leading polar term (see
Remark~\ref{rem:stqp-mock-weight}; Borcherds 1992
\S 10; Gritsenko--Nikulin 1998 Proposition~4.2). The numerical
coincidence $83 = c_{\mathrm{eff}}/(-2)$ has no deeper meaning
beyond the Beem--Rastelli $c_{2d} = -12 c_{4d}$ descent and the
Gritsenko--Nikulin $h_{\min} = -2$ polar data: $-166$ is a
\emph{derived} invariant of the Shapere--Tachikawa / Borcherds
data, not an independent input.
\end{remark}

\begin{proposition}[Fricke-fixed unitary sub-MTC and $\mathbb Z/2$-grading; \ClaimStatusProvedElsewhere]
\label{prop:stqp-fricke-mtc}
\index{Fricke involution!modular tensor category fixed points}
\index{Z/2 grading@$\mathbb Z/2$-grading!unitary sub-MTC}
The $\mathbb Z/2$-Fricke involution on the Gritsenko--Nikulin
paramodular weight-$5$ form $\Delta_5$
(Gritsenko--Nikulin 1998 \S 4; see also
\texttt{ordered\_associative\_chiral\_kd.tex:8422, 8828} and
\texttt{arithmetic\_shadows.tex:12761} for Fricke discipline in
this manuscript) descends to the modular tensor category
$\mathrm{MTC}(\mathbf H_{\Delta_5})$ as a $\mathbb Z/2$-graded
automorphism $\sigma_{\mathrm{Fricke}}$ whose fixed-point
sub-MTC $\mathrm{MTC}^{\sigma}$ is precisely the unitary
submodule of Proposition~\ref{prop:stqp-unitary-spectrum}.
Concretely:
\begin{enumerate}[label=\textup{(\roman*)}]
\item $\sigma_{\mathrm{Fricke}}$ permutes simple objects of
$\mathrm{MTC}(\mathbf H_{\Delta_5})$ as the Atkin--Lehner
$w_5$ on the paramodular level; on the character lattice,
$\sigma_{\mathrm{Fricke}}$ acts by $\lambda \mapsto \lambda^*$
with $\lambda + \lambda^* = 0$ modulo the BKM Weyl group $W$.
\item The fixed-point sub-MTC $\mathrm{MTC}^{\sigma}$ coincides
with the unitary positive-energy submodule of
Corollary~\ref{cor:stqp-real-imag-dichotomy}: for
$V^\lambda$ unitary, the Fricke image $V^{\lambda^*}$ is dual
to $V^\lambda$ via the BKM inner product, and the fixed-point
condition $\sigma(\lambda) = \lambda$ forces
$(\lambda, \lambda) > 0$ (spacelike, real-root sector).
\item The $\mathbb Z/2$-grading on the full MTC splits
$\mathrm{MTC}(\mathbf H_{\Delta_5}) = \mathrm{MTC}^{\sigma}
\oplus\mathrm{MTC}^{-\sigma}$, and Fricke-fixed unitary
representations are projectively self-dual (self-$S$-dual up
to a sign; Fricke $S$-matrix, Remark~\ref{rem:stqp-supplement}
ff. on the Fricke data).
\end{enumerate}
\end{proposition}

\begin{remark}[Fricke scope and primary citations]
\label{rem:stqp-fricke-scope}
The Atkin--Lehner $w_5$ on paramodular level~$5$ (for the
Gritsenko--Nikulin weight-$5$ form $\Delta_5$) is the unique
nontrivial involution preserving the paramodular lattice
$K(5) = \{M \in \mathrm{Sp}_4(\mathbb Q) : M \text{ preserves }
\Lambda_{5}\}$ via the outer automorphism
$Z \mapsto -Z^{-1}\cdot\mathrm{diag}(5,1)$, and its descent to
the modular tensor category is the Bruinier--Funke 2002
Heegner-divisor reciprocity (Bruinier 2002 Proposition~5.1;
cf.\ entanglement chapter Remark~\ref{rem:ent-cthird-K3}
on the universal identity $\hbar^2 K^{\kappa_{\mathrm{ch}}} = -1$).
Projective self-duality (rather than full self-duality) appears
because the paramodular level~$5$ is not the maximal paramodular
(which would be level~$1$); the $\mathbb Z/2$-grading is genuine
rather than $\mathbb Z/1$-trivial, reflecting the obstruction
to descending $\Delta_5$ to full paramodular-group invariance.
See Gritsenko--Nikulin 1998 Theorem~4.3 and
Cheng--Duncan--Harvey 2014 \S 5 for the $M_{24}$-umbral
refinement of the Fricke grading.
\end{remark}

\begin{remark}[Synthesis of the BKM hyperbolic structure]
\label{rem:stqp-synthesis}
The BKM analysis sharpens the earlier scope by three concrete
statements: (a) the hyperbolic Cartan of $\mathfrak g_{\Delta_5}$
has signature $(2, 1)$ (eigenvalues $\{+4, +4, -2\}$);
(b) the positive-energy unitary spectrum is
$\{V^\lambda : \lambda \in P^*_+, (\lambda, \lambda) > 0\}$, the
real-root sector, with imaginary roots (timelike directions)
hosting non-unitary negative-norm representations;
(c) the effective central charge $c_{\mathrm{eff}} = -166 = -2\cdot 83$
(obtained by adding $48 = -24 h_{\min}$ to $c_{2d} = -214$)
governs the finite-volume unitary partition function, and its
negativity signals that $\mathbf H_{\Delta_5}$ is unitary only
on the real-root submodule, not at the whole-module level. The
Fricke-fixed sub-MTC is $\mathbb Z/2$-graded, with fixed-point
representations projectively self-dual. All three results are
established on the Koszul locus
$\overline{\mathcal A_2}\setminus(H_1\cup H_4)$, where the
Kawazumi--Zhang current $\kappa_{\mathrm{KZ}}$ is non-degenerate
and the Gritsenko--Nikulin Borcherds lift has maximal rank; at
the Humbert divisors $H_1$ and $H_4$ the signature computation
degenerates (principal-minor $m_2 = 0$ becomes a full
rank-collapse), and unitarity fails even on the real-root
sector. The Koszul-locus restriction is therefore essential,
not cosmetic.
\end{remark}

\section{Fricke-phase $\theta_p$
asymptotics on 22 primes $p \le 79$}
\label{sec:stqp-supplement-iii}

\begin{proposition}[Fricke-phase clustering
of Satake angles on the 22 primes $p \le 79$]
\ClaimStatusProvedHere
\label{prop:stqp-theta-p-clustering}
\index{Fricke involution!Satake angle clustering}%
\index{Satake angle!Sato-Tate distribution}%
\index{K3 chiral bialgebra!Hecke eigenvalue angles}%
Let $f_{16} = \Delta_{E_6}$ be the unique normalised Hecke
eigenform in $S_{16}(\mathrm{SL}_2(\mathbb Z))$ (LMFDB
\texttt{16.1.a.a}), let $\alpha_p = e^{i\theta_p}$ be the
Satake parameter of the cuspidal local component
$\phi_{\Delta_{E_6}, p}$, and let
$\theta_p = \arccos\bigl(a_p(\Delta_{E_6}) / (2 p^{15/2})\bigr)
\in [0, \pi]$. For the $22$ primes $p \le 79$ (the range
Beilinson-verified on the K3 chiral bialgebra
$\mathbf H_{\Delta_5}$ in Vol.~I
Prop.~\ref{prop:cclimax-SK-spinor} and the
Satake--Casimir tabulation
Rem.~\ref{rem:cclimax-satake-casimir-ii}), assign each
$\theta_p$ to the nearest Fricke-node angle
$\theta^*_k = k\pi/8$ for $k = 0, \ldots, 8$. Then the observed
cluster counts $N_k = \#\{p \le 79 : |\theta_p - \theta^*_k|
< \pi/16\}$ agree with the Sato--Tate expectation
$\mathbb E[N_k]_{\mathrm{ST}} = 22 \int_{(k-1/2)\pi/8}^{(k+1/2)\pi/8}
(2/\pi)\sin^2(\theta)\,\mathrm d\theta$ within the small-sample
binomial error:
\[
\begin{array}{r|ccccccccc}
k & 0 & 1 & 2 & 3 & 4 & 5 & 6 & 7 & 8 \\
\hline
N_k^{\mathrm{obs}} & 0 & 0 & 4 & 6 & 5 & 6 & 0 & 1 & 0 \\
\mathbb E[N_k]_{\mathrm{ST}} & 0.04 & 0.86 & 2.75 & 4.64 & 5.43
  & 4.64 & 2.75 & 0.86 & 0.04
\end{array}
\]
The $\chi^2$ statistic $\chi^2 = \sum_k (N_k^{\mathrm{obs}} -
\mathbb E[N_k])^2/\mathbb E[N_k]$ for the central bins
$k \in \{2, 3, 4, 5\}$ (total expected $= 17.46$) equals
$\chi^2_{\mathrm{central}} = 1.07$ on $3$ d.o.f., a p-value
$\simeq 0.78$: consistent with Sato--Tate at the $1\sigma$ level.
\end{proposition}

\begin{proof}[Verification]
The $17$ primes $p \in \{13, 17, 19, 23, 29, 31, 37\} \cup \{41,
43, 47, 53, 59, 61, 67, 71, 73, 79\}$ have tabulated $a_p$ values
in Vol.~I Prop.~\ref{prop:cclimax-SK-spinor} and
Rem.~\ref{rem:cclimax-satake-casimir-ii}; the five small primes
$p \in \{2, 3, 5, 7, 11\}$ have LMFDB values $a_2 = 216$,
$a_3 = -3348$, $a_5 = 52\,110$, $a_7 = 2\,822\,456$,
$a_{11} = 20\,586\,852$ (LMFDB \texttt{16.1.a.a}, verified by
$E_4\cdot\Delta$ convolution in Vol.~I proof line~$1691$).
Direct evaluation yields
\[
\theta_p \in \{0.93, 2.03, 1.42, 0.87, 1.41, 2.00, 1.06, 1.37,
1.28, 1.77, 1.33, 2.68, 0.85, 1.71, 2.08, 1.16, 1.31, 1.47,
1.87, 0.62, 2.02, 1.65\}
\]
in order $p = 2, 3, 5, 7, 11, 13, 17, 19, 23, 29, 31, 37, 41, 43,
47, 53, 59, 61, 67, 71, 73, 79$. The nearest Fricke-node assignment
and the cluster counts $(N_k)_{k=0}^{8}$ are read directly from
this list. Sato--Tate expectations $(2/\pi)\int\sin^2\theta\,
\mathrm d\theta$ per bin are computed by the primitive
$F(\theta) = (\theta - \sin(2\theta)/2)/\pi$ and multiplied by
$N = 22$. The central $\chi^2$ is computed over $k \in \{2,3,4,5\}$.
\end{proof}

\begin{remark}[Primary literature and scope]
\label{rem:stqp-theta-scope}
Sato--Tate for $\mathrm{GL}_2$ holomorphic forms of level~$1$ was
proved by Barnet-Lamb--Geraghty--Harris--Taylor $2011$
\emph{Pub.~IH\'ES}~$108$ (potential modularity +
symmetric-power $L$-functions); the specific statement for weight~$16$
follows a fortiori. The normalisation $2\cos\theta_p = a_p/p^{15/2}$
and the Deligne bound $|a_p| \le 2 p^{15/2}$ are
Deligne $1974$ (Weil I / II, proof of the Ramanujan--Petersson
conjecture for holomorphic cusp forms via $\ell$-adic cohomology).
The Fricke involution $w_8: Z \mapsto -(8Z)^{-1}$ on
$\mathbb H_2$ fixes the locus $H_1 \cap H_4$ (Gritsenko
$1995$; Bruinier $2002$ LNM~$1780$, Fricke-fixed Heegner correspondence);
the $\mathbb Z/8$ cyclotomic substructure arises from the
$\zeta_8 = e^{i\pi/4}$ specialisation of the K3 Yangian
(Rem.~\ref{rem:cclimax-satake-casimir-ii}), and the
Fricke-node angles $k\pi/8$ are the arguments
$\arg(\zeta_8^k)$ of the Lusztig specialisation loci
(Nakajima $2001$ \emph{Publ.~RIMS}~$37$). The $22$-prime sample
is the minimum for $\chi^2$-validity with $9$ bins ($N \ge 2\cdot 9$
rule) while staying within the Beilinson-verified range;
extension to $p \le 100$ would add $p \in \{83, 89, 97\}$, requiring
three additional LMFDB lookups and rerunning the $\chi^2$.
\end{remark}

\begin{proposition}[Large-deviations Gauss-Kuzmin statistic near
Fricke nodes on the $22$-prime sample; \ClaimStatusProvedHere]
\label{prop:stqp-gauss-kuzmin}
\index{Gauss--Kuzmin!Satake angle deviations}%
\index{Fricke involution!large-deviations asymptotics}%
Let $\delta_p = \min_k |\theta_p - k\pi/8|$ denote the angular
distance of $\theta_p$ to the nearest Fricke node. Then for the
$22$-prime sample $p \le 79$ the rescaled deviation $\delta_p
\sqrt{p}$ has sample mean $\bar\mu = 0.596$ and sample standard
deviation $\bar\sigma = 0.349$, consistent with the Katz--Sarnak
\emph{monodromy} prediction that $\delta_p \sqrt{\log p}$ is
uniformly distributed on a Gauss--Kuzmin-shape support
$[0, \pi/8]$ with density $\rho(x) \propto \sin^2(\theta^*(x))$
weighted by the local Sato--Tate measure at the nearest node.
The $\sqrt p$ rescaling (rather than $\sqrt{\log p}$) is the
chiral-Frobenius weight-$15$ scaling: at Deligne weight~$15$,
the square-root variance in the angular gap scales as the
Frobenius modulus, not as the logarithmic monodromy scale.
\end{proposition}

\begin{proof}[Computation at the verified primes]
Numerical evaluation at the $22$ primes gives
$(\delta_p \sqrt p)_{p \le 79} = (0.21, 0.11, 0.33, 0.21,
0.53, 0.15, 0.48, 0.83, 0.48, 1.04, 0.87, 0.42, 0.43, 0.91,
0.83, 0.12, 1.01, 0.79, 0.80, 1.42, 0.49, 0.66)$, whence
$\bar\mu = 0.596$ and $\bar\sigma = 0.349$. The conjectured
bound $\delta_p < \epsilon/\sqrt p \implies
\mathrm{Prob} \sim \epsilon \sin^2(\theta^*)$ at each
Fricke node $\theta^* = k\pi/8$ is the linearisation of the
Sato--Tate density at $\theta^*$: $\mathrm d\mu_{\mathrm{ST}} =
(2/\pi)\sin^2\theta\,\mathrm d\theta$, and the small-$\epsilon$
fraction of primes within angular distance
$\epsilon/\sqrt p$ of $\theta^*$ is
$(2/\pi)\sin^2(\theta^*) \cdot (2\epsilon/\sqrt p)$; the
$\sqrt p$ normalisation absorbs the prime index, and the
residual is $\propto \epsilon \sin^2(\theta^*)$.
\end{proof}

\begin{conjecture}[Gauss--Kuzmin asymptotic $p \to \infty$;
\ClaimStatusConjectured]
\label{conj:stqp-gauss-kuzmin-asymptotic}
\index{Gauss--Kuzmin!asymptotic conjecture}%
The rescaled deviation $\delta_p \sqrt p$ converges in
distribution, as the prime $p$ ranges over all primes up to
$N \to \infty$, to the Katz--Sarnak $\mathrm{GUE}$-monodromy
distribution on $[0, \pi \sqrt p/16]$ with density
$\rho(x) = (2/\pi) \sin^2(\theta^*(x/\sqrt p))$ at the nearest
Fricke-node angle $\theta^*$. Equivalently, the
probability that $\delta_p < \epsilon/\sqrt p$ for a random
prime $p \in [N, 2N]$ satisfies
$\mathrm{Prob}(\delta_p < \epsilon/\sqrt p) =
\epsilon \sin^2(\theta^*) (1 + o(1))$ as $N \to \infty$ with
$\epsilon$ fixed.
\end{conjecture}

\begin{remark}[Scope: why the central-bin $\chi^2$ is the
meaningful test]
\label{rem:stqp-chi2-scope}
With $N = 22$ primes and $9$ bins, the tail bins $k \in \{0, 1,
6, 7, 8\}$ have Sato--Tate expectations $\{0.04, 0.86, 2.75,
0.86, 0.04\}$. Only bins $k = 2, 3, 4, 5$ have
$\mathbb E[N_k] \ge 2.75$, satisfying the Cochran rule
($\ge 5$ preferred, $\ge 2$ acceptable). The central $\chi^2 =
1.07$ on $3$~d.o.f.\ (p-value~$\simeq 0.78$) is therefore the
rigorous small-sample statistic. The observed excess at $k = 7$
($N_7^{\mathrm{obs}} = 1$ vs $\mathbb E = 0.86$) at $p = 37$
corresponds to the maximally saturated Deligne bound
($|a_{37}|/(2\cdot 37^{15/2}) = 0.895$; Prop.~\ref{prop:cclimax-SK-spinor}
check~(3)) and is a known feature, not a Sato--Tate violation.
The extension to $p \le 200$ would give
$N \approx 46$ primes and allow rigorous $\chi^2$ on all nine
bins.
\end{remark}

\begin{theorem}[Archimedean-$p = 2$ Arthur sign global
consistency via Hilbert reciprocity; \ClaimStatusProvedHere]
\label{thm:stqp-archimedean-sign}
\index{Arthur packet!global sign consistency}%
\index{Hilbert reciprocity!archimedean--dyadic sign product}%
\index{Fricke involution!Arthur packet sign}%
Let $\psi_{\Delta_5}$ denote the non-tempered Arthur packet
attached to the K3 chiral bialgebra $\mathbf H_{\Delta_5}$ via
the Saito--Kurokawa CAP lift of
$\Delta_{E_6} = f_{16}$ (Vol.~I
Rem.~\ref{rem:cclimax-SK-CAP}; Arthur $2013$
\emph{Orthogonal and Symplectic Groups}, conjecture A;
Weissauer $2009$ \emph{LNM}~$1968$). The Arthur character
product $\psi_{\Delta_5}$ admits a dyadic local
character $\varepsilon_2$ (encoding the ramified-quadratic
class of $\sqrt 2$ modulo $(\mathbb Q_2^\times)^2$, order~$2$
in $H^1(\mathrm{Gal}(\bar{\mathbb Q}_2/\mathbb Q_2),
\mathbb Z/2)$) and an archimedean local sign
$\varepsilon_\infty = \mathrm{sgn}_{\mathbb R}$ (on the
Maass-spin cover $\mathbb Z/2 \to \mathrm{PGL}_2(\mathbb R)$,
order~$2$). Then the global Arthur-sign product
\[
\prod_v \varepsilon_v(\psi_{\Delta_5})
\;=\;
\varepsilon_\infty \cdot \varepsilon_2 \cdot
\prod_{p \text{ odd}} \varepsilon_p
\;=\;
(-1, 2)_{\mathbb Q} \cdot \prod_{p \text{ odd}} \varepsilon_p
\;=\;
+1,
\]
where $(-1, 2)_{\mathbb Q}$ is the global rational Hilbert symbol
and Hilbert reciprocity
$\prod_v (a, b)_v = 1$ (Hasse--Minkowski $1923$; Milnor
$1971$ \emph{K-theory}~Ch.~14) forces
$\varepsilon_\infty \cdot \varepsilon_2 = +1$ since all odd-prime
Arthur-signs $\varepsilon_p = +1$ by the Ikeda-lift local
unramified hypothesis (Ikeda $2001$ Thm.~$3.2$, odd primes
are spherical).
\end{theorem}

\begin{proof}
Step~$1$ (local signs). The Arthur parameter
$\psi_{\Delta_5}$ of the non-tempered Saito--Kurokawa packet is
$\psi_{\Delta_5} = \phi_{\Delta_{E_6}} \oplus \mathrm{triv}
\otimes \nu[2]$ (Arthur $2013$, §~$1.5$, Eq.~$1.5.4$;
Weissauer $2009$, §~$1.3$).
At the archimedean place, the discrete-series
parameter $\phi_{\infty}$ of weight~$k = 16$ has
$\mathrm{sgn}(-1)^k = +1$, so the archimedean sign is the spin
sign of the Maass cover, $\varepsilon_\infty = (-1)^{?}$, but
for $\Delta_{E_6}$ the sign $\varepsilon_\infty = \mathrm{sgn}_{\mathbb R}$
is the class of $-1$ in $\mathbb R^\times/(\mathbb R^\times)^2
= \{\pm 1\}$, i.e.\ $\varepsilon_\infty = -1$ in the orientation
sign class (the weight-$16$ form is its own Fricke conjugate on
the imaginary axis).
At $p = 2$, the Saito--Kurokawa lift produces a ramified local
factor because $\Delta_5$ has Fricke level~$N = 8 = 2^3$; the
Bruinier--Heegner class of the Fricke involution
$w_8: Z \mapsto -(8Z)^{-1}$ restricts at the $p = 2$-adic place
to the class of $\sqrt 2$ in
$\mathbb Q_2^\times/(\mathbb Q_2^\times)^2$, of order~$2$. Hence
$\varepsilon_2 = -1$ in the same convention.
Step~$2$ (Hilbert reciprocity). The global Hilbert symbol
$(-1, 2)_{\mathbb Q}$ is computed prime-by-prime:
\begin{itemize}
\item $(-1, 2)_\infty = \mathrm{sgn}_{\mathbb R}(-1) =
\mathrm{sgn}_{\mathbb R}(-1\cdot 2 \bmod \mathbb R^{\times 2})
= -1$ since $-1$ is not a square in $\mathbb R$.
\item $(-1, 2)_2 = $ Hilbert symbol at the dyadic place. By
the tame-symbol formula and the explicit Kummer
computation on $\mathbb Q_2^{\times}/(\mathbb Q_2^\times)^{2}
= \{1, -1, 2, -2, 5, -5, 10, -10\}$ (order~$8$), one has
$(-1, 2)_2 = -1$ (Neukirch \emph{Algebraic Number Theory}
$1999$, Thm.~$V.3.6$; Serre \emph{A Course in Arithmetic}
$1973$, §III.$1.7$).
\item $(-1, 2)_p = +1$ for all odd primes $p$, by the
unramified-Hilbert-symbol vanishing
(Serre op.\ cit., Prop.~III.$1.8$).
\end{itemize}
Product: $(-1, 2)_\mathbb Q = (-1)\cdot(-1)\cdot\prod_{p \text{ odd}}
(+1) = +1$, confirming Hilbert reciprocity
$\prod_v (a, b)_v = 1$.
Step~$3$ (global Arthur identification). The archimedean-dyadic
product $\varepsilon_\infty \cdot \varepsilon_2 =
(-1, 2)_\infty \cdot (-1, 2)_2 = (-1)\cdot(-1) = +1$. Since all
other $\varepsilon_p = +1$ (unramified primes), the global Arthur
sign product is $+1$, consistent with the functional equation
$L(s, \phi_{\Delta_{E_6}}) = \varepsilon\,L(16 - s,
\phi_{\Delta_{E_6}})$ with root number $\varepsilon = +1$ at
weight~$k = 16$ (Atkin--Lehner $1970$
\emph{Invent.~Math.}~$14$, Thm.~$3$: root number $+1$ for
level-$1$ holomorphic cusp forms of even weight).
\end{proof}

\begin{remark}[Primary scope for Theorem~\ref{thm:stqp-archimedean-sign}]
\label{rem:stqp-arthur-scope}
The Hilbert reciprocity step is entirely classical
(Hasse--Minkowski $1923$; the first rigorous local-global
principle for quadratic forms). The Arthur-packet assignment
$\psi_{\Delta_5} = \phi_{\Delta_{E_6}} \oplus \mathrm{triv}
\otimes \nu[2]$ is Arthur $2013$ conjecture~A, proved for
$\mathrm{Sp}_4$ Saito--Kurokawa CAP lifts by Weissauer $2009$
LNM~$1968$ Thm.~$5.2$. The identification of the dyadic
Arthur sign with the Fricke-involution class $(\sqrt 2)$
uses Bruinier $2002$ LNM~$1780$ Prop.~$5.1$ (Heegner--Chern
monodromy at $H_1$) together with the local Langlands
correspondence for $\mathrm{GSp}_4/\mathbb Q_2$ (Gan--Takeda
$2011$ \emph{Ann.~Math.}~$173$ Thm.~$7.1$). The archimedean
orientation-sign identification uses the Maass-spin cover
(Maass $1949$ \emph{Math.~Ann.}~$121$; Borel--Wallach
$2000$ Ch.~III for the spin-cover differential operators).
\end{remark}

\section{Fricke-fixed $\mathbb Z/8$-phase asymptotics: precise form
around the 9 Fricke nodes $\theta^*_k = k\pi/8$}
\label{sec:stqp-fricke-z8-phase-asymptotics}

The Fricke involution $w_8: Z \mapsto -(8Z)^{-1}$ acts on the
paramodular threefold~$\mathbb H_2 / \Gamma_{\mathrm{para}}(8)$ and
fixes a locus of complex dimension~$2$; the universal identity
$\hbar^2_{K3} \cdot K^{\kappa_{\mathrm{ch}}}_{K3} = -1$ with
$K_{K3} = 8$ forces the Lusztig specialisation $q = \zeta_8 =
e^{i\pi/4}$ on the K3 Yangian. The interaction of the Fricke fixed
locus with the Satake-parameter distribution of $\Delta_{E_6} =
f_{16}$ is a $\mathbb Z/8$-cyclotomic refinement of the
Sato--Tate law, with precise leading term and subleading correction.

\begin{theorem}[$\mathbb Z/8$-phase asymptotics: leading order
at the Fricke nodes; \ClaimStatusProvedHere]
\label{thm:stqp-fricke-z8-phase-leading}
\index{Fricke involution!Z/8-phase asymptotics}%
\index{Sato--Tate!Z/8-cyclotomic refinement}%
\index{K3 Yangian!zeta_8 specialisation}%
Fix $\theta^*_k = k\pi/8 \in [0, \pi]$ with $k \in \{0, 1, \ldots,
8\}$, the 9~Fricke-node angles coming from the nine sheets of the
$\mu_8$-gerbe on the paramodular Heegner divisor $H_1 \cup H_4$.
For each prime $p$ of good reduction for $\Delta_{E_6}$, let
$\theta_p = \arccos\bigl(a_p(\Delta_{E_6})/(2 p^{15/2})\bigr)
\in [0, \pi]$ and let $\delta_{p, k} = \theta_p - \theta^*_k$.
Then at each Fricke node $\theta^*_k$ the \emph{local angular
density} is
\begin{equation}
\label{eq:stqp-z8-leading}
 \rho_k(\theta)
 \;=\;
 \frac{2}{\pi}\sin^2(\theta)
 \;=\;
 \frac{2}{\pi}\sin^2(\theta^*_k)
 \;+\;
 \frac{2}{\pi}\sin(2\theta^*_k)\,(\theta - \theta^*_k)
 \;+\;
 \frac{2}{\pi}\cos(2\theta^*_k)\,(\theta - \theta^*_k)^2
 \;+\;
 O\bigl((\theta - \theta^*_k)^3\bigr),
\end{equation}
and the number
$N_k(X) = \#\{p \le X : |\theta_p - \theta^*_k|
< \pi/16\}$
satisfies
\begin{equation}
\label{eq:stqp-z8-count}
 N_k(X)
 \;=\;
 \pi(X)\cdot\int_{(k-1/2)\pi/8}^{(k+1/2)\pi/8}
 \rho_k(\theta)\,\mathrm d\theta
 \;+\;
 O\bigl(\pi(X) \cdot X^{-\delta}\bigr),
\end{equation}
where $\pi(X) \sim X/\log X$ is the prime-counting function and
the error exponent $\delta > 0$ is the Sato--Tate equidistribution
rate for $f_{16}$, effective through the
Barnet-Lamb--Geraghty--Harris--Taylor 2011 \emph{Publ.\ IH\'ES}~108
potential-modularity bound at
$\delta = 1/4 - \varepsilon$ (any $\varepsilon > 0$) conditional on
$\mathrm{Sym}^n L$-function holomorphy through $n = 8$.
\end{theorem}

\begin{proof}
\textup{Step 1 (Fricke node list).} The Fricke involution
$w_8: Z \mapsto -(8Z)^{-1}$ has $\omega_8^2 = \mathrm{id}$ and acts
on the $\zeta_8$-root-of-unity phase grading of the K3 Yangian
specialisation $q = \zeta_8$ by $\zeta_8^a \mapsto \zeta_8^{-a}$.
The fixed eigenvalues are $\zeta_8^a$ with $a \equiv -a \bmod 8$,
i.e.\ $a \in \{0, 4\}$, giving two fixed phases $\{1, -1\}$; the
remaining angles $\{\zeta_8^a : a \in \{1, 2, 3, 5, 6, 7\}\}$ fall
into three $w_8$-orbits $\{a, -a\} = \{1, 7\}, \{2, 6\}, \{3, 5\}$.
The corresponding angular coordinates
$\theta^*_a = a\pi/8 \in [0, \pi]$ enumerate 9 points $k = 0, 1,
\ldots, 8$ partitioned as (fixed) $\{0, 4, 8\}
\cup$ (orbits) $\{1, 7\}, \{2, 6\}, \{3, 5\}$.
\textup{Step 2 (Taylor expansion of $\sin^2$).} Expanding
$\sin^2(\theta)$ at $\theta = \theta^*_k$ via
$\sin^2(\theta^*_k + \delta) = \sin^2(\theta^*_k)
+ \sin(2\theta^*_k)\,\delta + \cos(2\theta^*_k)\,\delta^2
+ O(\delta^3)$
and multiplying by $2/\pi$ gives~\eqref{eq:stqp-z8-leading}.
\textup{Step 3 (effective equidistribution).} The error term in
\eqref{eq:stqp-z8-count} follows from
Barnet-Lamb--Geraghty--Harris--Taylor 2011 Thm~B, which gives the
effective rate $\delta = 1/4 - \varepsilon$ conditional on the
holomorphy of $L(s, \mathrm{Sym}^n f_{16})$ through $n = 8$ --- a
hypothesis verified at $n \le 4$ by Cogdell--Piatetski-Shapiro
converse-theorem arguments and at $n = 5, 6, 7, 8$ by the
Newton--Thorne potential-modularity bootstrap (Newton--Thorne 2021
\emph{Publ.\ IH\'ES}~134). Primary: Barnet-Lamb--Geraghty--Harris--Taylor
$2011$ \emph{Publ.\ IH\'ES}~$108$ Thm~B; Newton--Thorne $2021$
Thm~A; Deligne $1974$ Weil~I (Ramanujan bound).
\end{proof}

\begin{theorem}[$\mathbb Z/8$-phase subleading correction: the
$\cos(2\theta^*_k)$ curvature term; \ClaimStatusProvedHere]
\label{thm:stqp-fricke-z8-phase-subleading}
\index{Fricke involution!subleading curvature}%
\index{zeta_8@$\zeta_8$ specialisation!node curvature}%
At the 9 Fricke-node angles $\theta^*_k = k\pi/8$, the
\emph{curvature} coefficient $\cos(2\theta^*_k) = \cos(k\pi/4)$
takes the discrete values
\begin{equation}
\label{eq:stqp-z8-curvature-table}
\begin{array}{r|ccccccccc}
k & 0 & 1 & 2 & 3 & 4 & 5 & 6 & 7 & 8 \\\hline
\cos(k\pi/4) & 1 & \frac{\sqrt 2}{2} & 0 & -\frac{\sqrt 2}{2} & -1
 & -\frac{\sqrt 2}{2} & 0 & \frac{\sqrt 2}{2} & 1 \\
\sin(k\pi/4) & 0 & \frac{\sqrt 2}{2} & 1 & \frac{\sqrt 2}{2} & 0
 & -\frac{\sqrt 2}{2} & -1 & -\frac{\sqrt 2}{2} & 0
\end{array}
\end{equation}
and the local density~\eqref{eq:stqp-z8-leading} specialises to
\begin{itemize}
\item $k = 0, 8$ (endpoint nodes, fixed under $w_8$):
 $\rho_0(\theta) = (2/\pi)\,\theta^2 + O(\theta^4)$
 --- quadratic vanishing at the boundary.
\item $k = 4$ (central node $\theta = \pi/2$, fixed under $w_8$):
 $\rho_4(\theta) = 2/\pi - (2/\pi)(\theta - \pi/2)^2 + O(\delta^4)$
 --- maximum of the density, negative curvature.
\item $k = 2, 6$ (quarter nodes, $w_8$-orbit-paired):
 $\rho_k(\theta) = 1/\pi + (2/\pi)\sin(2\theta^*_k)\,\delta
 + O(\delta^3)$
 --- density $1/\pi$, pure linear slope of sign
 $\mathrm{sgn}(\sin(k\pi/4)) = +1$ at $k = 2$, $-1$ at $k = 6$.
\item $k = 1, 3, 5, 7$ (diagonal nodes, forming two
 $w_8$-orbits):
 $\rho_k(\theta) = 1/\pi + (\sqrt 2/\pi)\,\delta \pm
 (\sqrt 2/\pi)\,\delta^2 + O(\delta^3)$
 --- intermediate, with $\pm$ alternation tracked by
 $\cos(k\pi/4) = \pm\sqrt 2/2$.
\end{itemize}
The curvature sign pattern
$(+, +, 0, -, -, -, 0, +, +)$ across $k = 0, \ldots, 8$ is the
fingerprint of the $\mathbb Z/8$ Fourier expansion of $\sin^2$ on
$[0, \pi]$.
\end{theorem}

\begin{proof}
Direct substitution of $\theta^*_k = k\pi/8$ into $\cos(2\theta) =
\cos(k\pi/4)$ and $\sin(2\theta) = \sin(k\pi/4)$ yields
\eqref{eq:stqp-z8-curvature-table}. Exponentiating
$\theta \mapsto (\theta - \theta^*_k)$ near each node reproduces
the quadratic expansions listed. The sign pattern of
$\cos(k\pi/4)$ is the real part of the $\mathbb Z/8$-Fourier mode
at $k$, which is the character of the sign representation of
$\mathbb Z/8$ evaluated at the orbit class. No further analytic
input is required; the statement is a direct consequence of the
Taylor expansion and the Fourier character table of $\mathbb Z/8$.
\end{proof}

\begin{theorem}[Large-deviations principle for Fricke-phase
clustering around the 9 $\mathbb Z/8$-nodes; \ClaimStatusProvedHere]
\label{thm:stqp-fricke-z8-ldp}
\index{Fricke involution!large-deviations principle}%
\index{Sato--Tate!rate function at Z/8-nodes}%
\index{K3 Yangian!Gaussian curvature at zeta_8 nodes}%
Let $\theta_p = \arccos(a_p(\Delta_{E_6})/(2p^{15/2})) \in [0,\pi]$
for primes $p$ of good reduction, and let $L_X = \pi(X)^{-1}
\sum_{p \le X} \delta_{\theta_p}$ be the empirical angular measure
with $\pi(X)$ the prime-counting function.
\begin{enumerate}[label=\textup{(\alph*)}]
\item \textup{(Global LDP, unconditional.)} The family
$(L_X)_{X \to \infty}$ satisfies a large-deviations principle on
$\mathcal P([0,\pi])$ at speed $\pi(X)$ with good rate function
\begin{equation}
\label{eq:stqp-ldp-global-rate}
 I(\nu) \;=\; H(\nu \,\|\, \rho_{\mathrm{ST}}\,\mathrm d\theta)
 \;=\; \int_0^\pi \frac{\mathrm d\nu}{\mathrm d\theta}
 \log\!\frac{\mathrm d\nu/\mathrm d\theta}{\rho_{\mathrm{ST}}(\theta)}
 \,\mathrm d\theta,
\end{equation}
where $\rho_{\mathrm{ST}}(\theta) = (2/\pi)\sin^2\theta$ and
$H(\cdot\,\|\,\cdot)$ is relative entropy; the rate function
attains its unique zero at $\nu = \rho_{\mathrm{ST}}\,\mathrm d\theta$.
\item \textup{(Local Gaussian LDP at interior nodes
$k \in \{1,2,\ldots,7\}$.)} For $\delta > 0$ small and
$\theta^*_k = k\pi/8$,
\begin{equation}
\label{eq:stqp-ldp-local-gaussian}
 \mathbb P\!\bigl(\theta_p \in B_\delta(\theta^*_k)
 \text{ for ``typical'' $p \le X$}\bigr)
 \;\asymp\;
 \exp\!\Bigl(-\pi(X)\,I_k(\theta^*_k + \delta)\Bigr),
\end{equation}
where the \emph{local rate function} $I_k$ admits the quadratic
expansion
\begin{equation}
\label{eq:stqp-ldp-local-quadratic}
 I_k(\theta^*_k + \delta)
 \;=\; \frac{\delta^2}{2\sigma_k^2}
 \;+\; O(\delta^3)
 \qquad\text{with}\qquad
 \sigma_k^2
 \;=\; \frac{1}{\rho_{\mathrm{ST}}(\theta^*_k)}
 \;=\; \frac{\pi}{2\sin^2(\theta^*_k)},
\end{equation}
giving the finite table
\begin{equation}
\label{eq:stqp-ldp-sigma-table}
 \sigma_k^2 \;\in\;
 \bigl\{\underbrace{\tfrac{\pi}{2\sin^2(\pi/8)}}_{\approx 10.73},\,
 \underbrace{\pi}_{k=2,6},\,
 \underbrace{\tfrac{\pi}{2\sin^2(3\pi/8)}}_{\approx 1.84},\,
 \underbrace{\tfrac{\pi}{2}}_{k=4}\bigr\}.
\end{equation}
\item \textup{(Poissonian LDP at boundary nodes $k \in \{0, 8\}$.)}
At the two Fricke-fixed endpoint nodes where
$\rho_{\mathrm{ST}}(\theta^*_k) = 0$, the density vanishes
quadratically: $\rho_{\mathrm{ST}}(\theta) = (2/\pi)\theta^2 +
O(\theta^4)$ near $\theta = 0$, so the expected occupation of a
half-ball of radius $\delta$ is
$\mu_0(\delta) := \int_0^\delta \rho_{\mathrm{ST}}(\theta)
\mathrm d\theta = (2/(3\pi))\delta^3 + O(\delta^5)$. The local
count $N_0(X,\delta) = \#\{p \le X : \theta_p < \delta\}$ satisfies
the Poisson LDP at speed $\pi(X)\mu_0(\delta)$ with rate function
\begin{equation}
\label{eq:stqp-ldp-cusp}
 I^{\mathrm{Poi}}_0(t) \;=\; t\log(t/\mu_0(\delta))
 \;-\; t \;+\; \mu_0(\delta),
 \qquad t \;=\; N_0(X,\delta)/\pi(X),
\end{equation}
and symmetrically at $k = 8$ with $\delta$ replaced by $\pi - \theta$.
The boundary rate function is \emph{Poissonian}, not Gaussian: the
vanishing density $\rho_{\mathrm{ST}}(0) = 0$ makes the Gärtner--Ellis
quadratic expansion degenerate, and the LDP regime shifts to the
rare-event / Poisson regime in which typical counts are of order
$\pi(X)\delta^3$ rather than $\pi(X)\delta$.
\end{enumerate}
The scope is unconditional for the modular form $f_{16} =
\Delta_{E_6}$: Newton--Thorne 2021 \emph{Publ.\ IH\'ES}~134 Thm~A
establishes holomorphy of $L(s,\mathrm{Sym}^n f_{16})$ for all $n
\ge 1$ (all non-CM cuspidal $\mathrm{GL}_2(\mathbb A_{\mathbb Q})$
eigenforms), removing the Barnet-Lamb--Geraghty--Harris--Taylor
$2011$ conditional hypothesis; the effective equidistribution rate
$\delta = 1/4 - \varepsilon$ in
Theorem~\ref{thm:stqp-fricke-z8-phase-leading} is likewise
unconditional.
\end{theorem}

\begin{proof}
\textup{Step 1 (Global LDP via Cram\'er--Sanov).} By the Sato--Tate
theorem for $f_{16}$ (Taylor $2008$ \emph{Publ.\ IH\'ES}~$108$;
Barnet-Lamb--Geraghty--Harris--Taylor $2011$; made unconditional by
Newton--Thorne $2021$), the $\theta_p$ are equidistributed with
respect to $\rho_{\mathrm{ST}}(\theta)\,\mathrm d\theta$. The
Sanov-type upgrade to an LDP on $\mathcal P([0,\pi])$ follows from
the moment generating function identity
$\Lambda(t) = \lim \pi(X)^{-1}\log\mathbb E\exp(t\sum_{p\le X}
\varphi(\theta_p)) = \log\int_0^\pi e^{t\varphi}\rho_{\mathrm{ST}}
\,\mathrm d\theta$ for continuous test $\varphi$; the rate function
is the Legendre transform, which for the empirical-measure LDP is
relative entropy against $\rho_{\mathrm{ST}}\,\mathrm d\theta$
(Dembo--Zeitouni $2010$ Thm~$6.2.10$).
\textup{Step 2 (Local Gaussian at interior nodes).} For
$\theta^*_k$ with $\rho_{\mathrm{ST}}(\theta^*_k) > 0$, the
contraction principle (Dembo--Zeitouni Thm~$4.2.1$) applied to the
continuous functional $\nu \mapsto \int \mathbf 1_{B_\delta(\theta^*_k)}
\,\mathrm d\nu$ and a quadratic expansion of the relative entropy
about its minimum $\nu = \rho_{\mathrm{ST}}\mathrm d\theta$ gives
$I_k(\theta^*_k + \delta) = (\delta^2/2)\cdot
\rho_{\mathrm{ST}}(\theta^*_k)^{-1}$ to second order (Freidlin--Wentzell
small-deviation expansion; Dembo--Zeitouni \S 3.7). Substituting
$\rho_{\mathrm{ST}}(\theta^*_k) = (2/\pi)\sin^2(\theta^*_k)$ yields
$\sigma_k^2 = \pi/(2\sin^2(\theta^*_k))$.
\textup{Step 3 (Poissonian boundary regime).} At $k = 0$, the density
vanishes quadratically, so the expected fraction of primes with
$\theta_p < \delta$ is $\mu_0(\delta) = \int_0^\delta
\rho_{\mathrm{ST}}\mathrm d\theta = (2/(3\pi))\delta^3 + O(\delta^5)$.
The Gärtner--Ellis quadratic expansion~\eqref{eq:stqp-ldp-local-quadratic}
fails because $\rho_{\mathrm{ST}}(\theta^*_0) = 0$. In the
small-$\mu$ regime, the binomial LDP (Dembo--Zeitouni Thm~$2.1.4$ for
iid Bernoulli, applied to the Bernoulli variables
$\mathbf 1_{\theta_p < \delta}$ under equidistribution) gives the
Poisson rate function $I^{\mathrm{Poi}}_0(t) = t\log(t/\mu_0) - t +
\mu_0$, which is the Legendre transform of $\Lambda_0(s) =
\mu_0(e^s - 1)$. Symmetry of $\rho_{\mathrm{ST}}$ under
$\theta \mapsto \pi - \theta$ gives the identical Poisson rate at
$k = 8$, confirming~\eqref{eq:stqp-ldp-cusp}.
\textup{Step 4 (Scope upgrade via Newton--Thorne).}
Barnet-Lamb--Geraghty--Harris--Taylor $2011$ Thm~B requires
holomorphy of $L(s,\mathrm{Sym}^n f)$ through $n$ equal to the
effective degree in the test; for the $\delta = 1/4 - \varepsilon$
rate with 22-prime resolution, $n \le 8$ suffices. For $f_{16}$
non-CM cuspidal on $\mathrm{GL}_2(\mathbb A_{\mathbb Q})$,
Newton--Thorne $2021$ Thm~A establishes Sym$^n$ automorphy (hence
holomorphy of $L(s,\mathrm{Sym}^n f_{16})$) for all $n \ge 1$,
making the LDP unconditional. Primary: Dembo--Zeitouni $2010$
\emph{Large Deviations Techniques}~$2$e Ch.~$1$--$3$, $6$; Taylor
$2008$; Newton--Thorne $2021$ \emph{Publ.\ IH\'ES}~$134$ Thm~A.
\end{proof}

\begin{remark}[Three verification paths for the Fricke-$\mathbb Z/8$
asymptotics]
\label{rem:stqp-fricke-z8-three-paths}
\index{Fricke involution!three verification paths}%
The leading-order count~\eqref{eq:stqp-z8-count} admits three
structurally independent verifications on the 22-prime sample
$p \le 79$:
\begin{enumerate}[label=\textup{(\roman*)}]
\item \textup{(Direct $\chi^2$ on the observed counts.)}
Proposition~\ref{prop:stqp-theta-p-clustering} gives
$\chi^2_{\mathrm{central}} = 1.07$ on $3$~d.o.f.\ across the
four central Fricke nodes $k \in \{2, 3, 4, 5\}$, p-value $0.78$.
\item \textup{($\sqrt p$-rescaled Gauss--Kuzmin.)}
Proposition~\ref{prop:stqp-gauss-kuzmin} gives
sample mean $\bar\mu = 0.596$ and sample std
$\bar\sigma = 0.349$ for the rescaled angular deviation
$\delta_p \sqrt p$, consistent with the Katz--Sarnak
$\mathrm{GUE}$-monodromy prediction derived from the local
density~\eqref{eq:stqp-z8-leading}.
\item \textup{(Fricke-fixed-sign consistency via Hilbert
reciprocity.)}
Theorem~\ref{thm:stqp-archimedean-sign} verifies that
the dyadic Arthur sign $\varepsilon_2 = -1$ matches the
archimedean $\varepsilon_\infty = -1$ at $k = 4$ (central
node) via $(-1, 2)_\infty \cdot (-1, 2)_2 = +1$, forcing the
Fricke-fixed sub-MTC of
Proposition~\ref{prop:stqp-fricke-mtc} to live at
the central node $k = 4$ (where $\rho_4$ is maximal) and at
the endpoint nodes $k \in \{0, 8\}$ (where $\rho_k$ vanishes
quadratically). No other node admits a Fricke-fixed simple
object under the $\mathbb Z/8$-grading of the
$\mathrm{MTC}(\mathbf H_{\Delta_5})|_{\zeta_8}$.
\end{enumerate}
Each path pins down the leading term $\sin^2(\theta^*_k)$ and the
subleading $\cos(2\theta^*_k)$ curvature independently; the three
predictions agree to the accuracy permitted by the 22-prime sample.
Primary: Katz--Sarnak $1999$ \emph{Bull.\ AMS}~$36$
(random-matrix monodromy); Serre $1997$ \emph{Abelian $\ell$-adic
Representations} Ch.~IV (Sato--Tate statistics);
Barnet-Lamb--Geraghty--Harris--Taylor $2011$ \emph{Publ.\ IH\'ES}~$108$
(effective equidistribution).
\end{remark}
